% 第 I 讲第 9--10 节及参考文献；原 PDF 第 86 页第 9 节起至第 112 页末。
\SourcePage{86}
\SGAsection{9}{可达函子、基数滤过与小生成子范畴的构造}
本节比本讲其他各节更具技术性，在本讨论班中只用于第 IV 讲第 9 节和第 VI 讲第 4 节，而这两节又不在本讨论班的其他地方使用。因此，应当略过本节，至少初读时应当如此！

\SourcePage{87}
\textbf{9.0。}本节所考虑的所有\Term{category}{范畴}都假定为 $\mathcal U$-\Term{category}{范畴}。除了我们将使用的小\Term{index-category}{指标范畴} $I,J,\ldots$ 以外，以下论述主要适用于对小\Term{filtered-inductive-limit}{滤过归纳极限}\Term{stable}{稳定}的“大”\Term{category}{范畴} $E,F,\ldots$。不过，通常只须满足一个稍弱的条件就够了（下文 9.1 中的条件 L）。本节所考虑的所有\Term{cardinal}{基数}都假定属于 $\mathcal U$。

按照 P.~Deligne 的一个建议，我们将研究大\Term{category}{范畴}之间的\Term{functor}{函子} $f:E\to F$ 与某些类型的\Term{filtered-inductive-limit}{滤过归纳极限}交换的条件。这个条件具有显著的\Term{stable}{稳定}性，通常遇到的最重要的\Term{functor}{函子}都满足它。对于本讨论班，我们所着眼的应用是 9.13.3、9.13.4（用于第 VI 讲第 4 节），尤其是 9.25：它在一个非\Term{trivial-bundle}{平凡}的情形中给出\Term{fiber-category}{纤维范畴}的\Term{section-category}{截面范畴}中小\Term{generating-family}{生成族}的存在性；这个结果将用于第 IV 讲 9.16。

\begin{SGAresult}{定义}{9.1}
a) 设 $I$ 是一个\Term{preordered-set}{预序集}，$\pi$ 是一个\Term{cardinal}{基数}。若 $I$ 是\Term{filtered-category}{滤过的}，且 $I$ 的每个\Term{cardinal}{基数} $\leq\pi$ 的\Term{subset}{子集}在 $I$ 中都有一个上界，就称 $I$ 相对于 $\pi$ \Term{large-relative-to-cardinal}{足够大}。

b) 设 $E$ 是一个\Term{category}{范畴}。若 $\pi$ 是一个\Term{cardinal}{基数}，当对每个相对于 $\pi$ \Term{large-relative-to-cardinal}{足够大} 的小\Term{ordered-set}{有序集} $I$，$E$ 都对 $I$ 型\Term{inductive-limit}{归纳极限}\Term{stable}{稳定}时，就称 $E$ 满足条件 $\mathrm L_\pi$。若存在\Term{cardinal}{基数} $\pi\in\mathcal U$，使得 $E$ 满足条件 $\mathrm L_\pi$，就称 $E$ 满足条件 $\mathrm L$。
\end{SGAresult}

\textbf{9.1.1。}当 9.1 a) 中的 $\pi\geq2$ 时，所述第二个条件本身已经蕴涵 $I$ 是\Term{filtered-category}{滤过的}；若 $\pi$ \Term{finite}{有限}，则 $I$ 相对于 $\pi$ \Term{large-relative-to-cardinal}{足够大} 仅仅意味着 $I$ 是\Term{filtered-category}{滤过的}。以下我们基本上只关心 $\pi$ 为无限\Term{cardinal}{基数}的情形。注意，若 $\pi,\pi'$ 是两个\Term{cardinal}{基数}，且 $\pi'\geq\pi$，则 $I$ 相对于 $\pi'$ \Term{large-relative-to-cardinal}{足够大} 显然蕴涵 $I$ 相对于 $\pi$ \Term{large-relative-to-cardinal}{足够大}。

\textbf{9.1.2。}如 9.0 中所说，条件 $\mathrm L_\pi$、$\mathrm L$ 应当看作更强的条件“对小\Term{filtered-inductive-limit}{滤过归纳极限}\Term{stable}{稳定}”的技术性变体。显然，若 $\pi,\pi'$ 是满足 $\pi'\geq\pi$ 的\Term{cardinal}{基数}，则条件 $\mathrm L_\pi$ 蕴涵条件 $\mathrm L_{\pi'}$。

\begin{SGAresult}{定义}{9.2}
设 $f:E\to F$ 是一个\Term{functor}{函子}。若 $\Pi$ 是一个\Term{cardinal}{基数}，当 $E$ 满足 $\mathrm L_\Pi$（9.1），且对每个相对于 $\Pi$ \Term{large-relative-to-cardinal}{足够大} 的\Term{ordered-set}{有序集} $I\in\mathcal U$ 及 $E$ 中每个 $I$ 型\Term{inductive-system}{归纳系统} $(X_i)_{i\in I}$，\Term{canonical}{典范}\Term{morphism}{态射}
\[
\varinjlim f(X_i)\longrightarrow f(\varinjlim X_i)
\]
是\Term{monomorphism}{单态射}（分别为\Term{isomorphism}{同构}）时，就称 $f$ 是 $\Pi$-\Term{preaccessible}{预可达}的（分别为 $\Pi$-\Term{accessible}{可达}的）。若存在\Term{cardinal}{基数} $\Pi\in\mathcal U$，使得 $f$ 是 $\Pi$-\Term{preaccessible}{预可达}的（分别为 $\Pi$-\Term{accessible}{可达}的），就称 $f$ 是\Term{preaccessible}{预可达}的（分别为\Term{accessible}{可达}的）（相对于\Term{universe}{宇宙} $\mathcal U$）。
\end{SGAresult}

\SourcePage{88}
从 $E$ 到 $F$ 的 $\Pi$-\Term{accessible-functor}{可达函子}（分别为\Term{accessible-functor}{可达函子}）所成的\Term{category}{范畴}记作 $\mathscr{H}om(E,F)_\Pi$，分别记作 $\mathscr{H}om(E,F)_{\mathrm{acc}}$。

\textbf{9.2.1。}显然，与小\Term{filtered-inductive-limit}{滤过归纳极限}交换的\Term{functor}{函子}（例如具有\Term{right-adjoint}{右伴随}的\Term{functor}{函子}），对于每个\Term{cardinal}{基数} $\Pi\geq2$ 都是 $\Pi$-\Term{accessible}{可达}的。

\begin{SGAresult}{定义}{9.3}
设 $E$ 是一个\Term{category}{范畴}，$X$ 是 $E$ 的一个\Term{object}{对象}，
\[
h_X^\circ:E\longrightarrow(\mathcal U\text{-}\mathrm{Ens})
\]
是它所表示的\Term{covariant-functor}{协变函子}，$\Pi\in\mathcal U$ 是一个\Term{cardinal}{基数}。若\Term{functor}{函子} $h_X^\circ$ 是 $\Pi$-\Term{preaccessible}{预可达}的（分别为 $\Pi$-\Term{accessible}{可达}的），就称 $X$ 是 $E$ 的一个 $\Pi$-\Term{preaccessible-object}{预可达对象}（分别为 $\Pi$-\Term{accessible-object}{可达对象}）；若存在\Term{cardinal}{基数} $\Pi\in\mathcal U$，使得 $X$ 是 $E$ 的 $\Pi$-\Term{preaccessible-object}{预可达对象}（分别为 $\Pi$-\Term{accessible-object}{可达对象}），就称 $X$ 是\Term{preaccessible}{预可达}的（分别为\Term{accessible}{可达}的）。
\end{SGAresult}

\textbf{9.3.1。}以 $E_\Pi$ 表示 $E$ 中由 $\Pi$-\Term{accessible-object}{可达对象}组成的\Term{full-subcategory}{满子范畴}。

将定义 9.3 用于形如 $\mathscr{H}om(C,F)$ 的\Term{category}{范畴}，并将 $\mathscr{H}om(C,F)$ 的\Term{object}{对象}解释为\Term{functor}{函子}时，这里引入的术语与 9.2 中引入的类似术语，先验地说有歧义；不过，似乎并没有严重混淆的危险。

\begin{SGAresult}{定义}{9.4}
设 $E$ 是一个\Term{category}{范畴}，$\pi\in\mathcal U$ 是一个\Term{cardinal}{基数}。若 $E$ 中存在一个小\Term{full-subcategory}{满子范畴} $C$，它是生成的（7.1），且其\Term{object}{对象}都是 $\pi$-\Term{preaccessible}{预可达}的（分别为 $\pi$-\Term{accessible}{可达}的）（9.3），就称 $E$ 是一个 $\pi$-\Term{preaccessible-category}{预可达范畴}（分别为 $\pi$-\Term{accessible-category}{可达范畴}）。若存在\Term{cardinal}{基数} $\pi\in\mathcal U$，使得 $E$ 是 $\pi$-\Term{preaccessible}{预可达}的（分别为 $\pi$-\Term{accessible}{可达}的），就称 $E$ 是一个\Term{preaccessible-category}{预可达范畴}（分别为\Term{accessible-category}{可达范畴}）。
\end{SGAresult}

重要的例子见下文 9.11.3。

\begin{SGAresult}{命题}{9.5}
设 $f:E\to F$ 是\Term{category}{范畴}之间的一个\Term{functor}{函子}，$F$ \Term{accessible}{可达}，且 $E$ 的每个\Term{object}{对象}都\Term{accessible}{可达}。若 $f$ 有\Term{left-adjoint}{左伴随}，则 $f$ \Term{accessible}{可达}。
\end{SGAresult}
\begin{proof}
由假设，$F$ 有一个\Term{accessible}{可达}的\Term{representable-functor}{可表示函子} $F\to(\mathcal U\text{-}\mathrm{Ens})$ 的小\Term{conservative-family}{保守族}，所以立即归结为证明 $f$ 与上述每个\Term{functor}{函子}的复合都是\Term{accessible}{可达}的。但因为 $f$ 有\Term{left-adjoint}{左伴随}，这些复合都是\Term{representable-functor}{可表示函子} $E\to(\mathcal U\text{-}\mathrm{Ens})$，故由关于 $E$ 的假设，它们都是\Term{accessible}{可达}的。
\end{proof}

\begin{SGAresult}{命题}{9.6}
设 $E,F$ 是两个\Term{category}{范畴}，$\pi\in\mathcal U$ 是一个\Term{cardinal}{基数}，并考虑 $\mathscr{H}om(E,F)$ 中由从 $E$ 到 $F$ 的 $\Pi$-\Term{accessible-functor}{可达函子}（9.2）组成的\Term{full-subcategory}{满子范畴} $\mathscr{H}om(E,F)_\Pi$。

(i) 这个\Term{subcategory}{子范畴}对 $F$ 中\Term{representable}{可表示}的任意类型的\Term{inductive-limit}{归纳极限}\Term{stable}{稳定}。
\SourcePage{89}
(ii) 假定 $F$ 是 $\pi$-\Term{accessible}{可达}的（9.4），且 $J$ 是一个满足 $\operatorname{card}\operatorname{Fl}J\leq\pi$ 的\Term{category}{范畴}，$J$ 型\Term{projective-limit}{投射极限}在 $F$ 中\Term{representable}{可表示}，因而在 $\mathscr{H}om(E,F)$ 中也\Term{representable}{可表示}。则\Term{subcategory}{子范畴} $\mathscr{H}om(E,F)_\Pi$ 对 $J$ 型\Term{projective-limit}{投射极限}\Term{stable}{稳定}。
\end{SGAresult}

\begin{SGAresult}{推论}{9.7}
设 $E,F$ 是两个\Term{category}{范畴}，且 $F$ \Term{accessible}{可达}（9.4）。则 $\mathscr{H}om(E,F)$ 中由\Term{accessible-functor}{可达函子}组成的\Term{full-subcategory}{满子范畴} $\mathscr{H}om(E,F)_{\mathrm{acc}}$，对任意以小\Term{index-category}{指标范畴} $J$ 为指标、在 $F$ 中（因而也在 $\mathscr{H}om(E,F)$ 中）\Term{representable}{可表示}的\Term{inductive-limit}{归纳极限}或\Term{projective-limit}{投射极限}\Term{stable}{稳定}。
\end{SGAresult}

\begin{proof}[9.6 的证明]
断言 (i) 直接来自\Term{inductive-limit}{归纳极限}\Term{functor}{函子}与任意\Term{inductive-limit}{归纳极限}交换。对于 (ii)，设 $(f_j)_{j\in J}$ 是从 $E$ 到 $F$ 的 $\Pi$-\Term{accessible-functor}{可达函子}的一个\Term{projective-system}{投射系统}，$f=\varprojlim f_j$ 是其\Term{projective-limit}{投射极限}，它“逐变量”计算。证明 $f$ 是 $\Pi$-\Term{accessible}{可达}的，即对每个相对于 $\Pi$ \Term{large-relative-to-cardinal}{足够大} 的\Term{ordered-set}{有序集} $I\in\mathcal U$ 和 $E$ 中每个\Term{inductive-system}{归纳系统} $(X_i)_{i\in I}$，\Term{canonical}{典范}\Term{morphism}{态射}
\[
\varinjlim_i\bigl((\varprojlim_j f_j)(X_i)\bigr)
\longrightarrow (\varprojlim_j f_j)(\varinjlim_i X_i)
\]
是\Term{isomorphism}{同构}。对\Term{functor}{函子} $\varprojlim_j f_j$ 的“逐变量”计算，使我们可以将上述\Term{canonical}{典范}\Term{morphism}{态射}认同于\Term{bifunctor}{双函子}
\[
(i,j)\longmapsto f_j(X_i):I\times J\longrightarrow F
\]
所关联的\Term{canonical}{典范}\Term{morphism}{态射}
\[
\varinjlim_i\varprojlim_j f_j(X_i)
\longrightarrow\varprojlim_j\varinjlim_i f_j(X_i).
\]
因此 9.6 是下述更一般断言的结果。
\end{proof}

\begin{SGAresult}{推论}{9.8}
设 $F$ 是一个 $\Pi$-\Term{accessible-category}{可达范畴}，$I$ 是一个相对于 $\Pi$ \Term{large-relative-to-cardinal}{足够大} 的\Term{ordered-set}{有序集}，$J$ 是一个满足 $\operatorname{card}\operatorname{Fl}\leq\Pi$ 的\Term{small-category}{小范畴}，且 $J$ 型\Term{projective-limit}{投射极限}在 $F$ 中\Term{representable}{可表示}，$h:I\times J\to F$ 是一个\Term{functor}{函子}。则\Term{canonical}{典范}\Term{morphism}{态射}
\[
\varinjlim_i\varprojlim_j h(i,j)
\longrightarrow\varprojlim_j\varinjlim_i h(i,j)\tag{9.8.1}
\]
是\Term{isomorphism}{同构}。换言之，\Term{functor}{函子}
\[
\varinjlim_i:\mathscr{H}om(I,F)\longrightarrow F\tag{9.8.2}
\]
与 $J$ 型\Term{projective-limit}{投射极限}交换（对每个满足 $\operatorname{card}\operatorname{Fl}(J)\leq\Pi$、且 $J$ 型\Term{projective-limit}{投射极限}在 $F$ 中\Term{representable}{可表示}的\Term{small-category}{小范畴} $J$）。也可以说，对每个如上的 $J$，\Term{functor}{函子}
\[
\varinjlim_j:\mathscr{H}om(J,F)\longrightarrow F\tag{9.8.3}
\]
是 $\Pi$-\Term{accessible}{可达}的。
\end{SGAresult}

\SourcePage{90}
\begin{proof}
由假设，$F$ 有一个 $\Pi$-\Term{accessible}{可达}的\Term{covariant}{协变}\Term{representable-functor}{可表示函子} $g:F\to(\mathcal U\text{-}\mathrm{Ens})$ 的\Term{conservative-family}{保守族}，所以立即看出只须在 $F=\mathcal U\text{-}\mathrm{Ens}$ 的情形中证明 9.8。我们将以 (9.8.2) 与 $J$ 型\Term{projective-limit}{投射极限}交换的形式来证明，其中 $\operatorname{card}\operatorname{Fl}J\leq\Pi$。由于 $I$ 是\Term{filtered-category}{滤过的}，我们知道 (9.8.2) 与\Term{finite}{有限}\Term{projective-limit}{投射极限}交换（2.8）。因此，由一个标准论证（参见 2.3），归结为证明它与以满足 $\operatorname{card}J\leq\Pi$ 的\Term{set}{集合} $J$ 为指标的积交换。这又归结为在 $J$ \Term{discrete}{离散}时证明 (9.8.1) 的\Term{bijection}{双射}性。

先证明\Term{injection}{单射}性。考虑左端的两个元素 $a,b$；它们来自某个适当的 $i_\circ\in I$ 所\Term{correspondence}{对应}的 $\prod_jh(i_\circ,j)$，设为 $\prod_j a(i_\circ,j)$ 和 $\prod_j b(i_\circ,j)$。假定它们在右端所定义的 $\prod_jh(\infty,j)$ 的元素相等，即对每个 $j$，存在 $i(j)\geq i_\circ$，使得 $a(i_\circ,j)$ 与 $b(i_\circ,j)$ 在 $h(i,j)$ 中有相同的像。由于 $I$ 相对于 $\operatorname{card}J$ \Term{large-relative-to-cardinal}{足够大}，所有 $i(j)$ 在 $I$ 中有一个共同上界 $i_1$。这意味着所考虑的 $\prod_jh(i_\circ,j)$ 的两个元素在 $\prod_jh(i_1,j)$ 中有相同的像，因而定义 (9.8.1) 左端的同一个元素，\Term{injection}{单射}性得证。

对于\Term{surjection}{满射}性，设 $\prod_j a(\infty,j)$ 是右端的一个元素。对每个 $j\in J$，$a(\infty,j)$ 都来自某个适当的 $i(j)\in I$ 所\Term{correspondence}{对应}的 $h(i(j),j)$ 中的一个元素。如前，可以在 $I$ 中找到所有 $i(j)$ 的共同上界 $i$，故 $\prod_j a(\infty,j)$ 来自 $\prod_jh(i,j)$ 中的一个元素 $\prod_j a(i,j)$，因而属于 (9.8.1) 的像。这就完成了证明。
\end{proof}

\begin{SGAresult}{推论}{9.9}
设 $F$ 是一个 $\Pi$-\Term{accessible-category}{可达范畴}（分别为\Term{accessible-category}{可达范畴}）。则对于每个满足 $\operatorname{card}\operatorname{Fl}J\leq\Pi$ 的\Term{category}{范畴} $J$（分别为每个\Term{small-category}{小范畴} $J$），和每个在 $F$ 中的\Term{inductive-limit}{归纳极限}\Term{representable}{可表示}的\Term{functor}{函子} $(X_j)_{j\in J}:J\to F$，若所有 $X_j$ 都是 $\Pi$-\Term{accessible}{可达}的（分别为\Term{accessible}{可达}的），则 $\varprojlim X_j$ 也是如此。
\end{SGAresult}
\begin{proof}
事实上，由 $\varinjlim X_j$ 所表示的\Term{functor}{函子} $F\to(\mathcal U\text{-}\mathrm{Ens})$ 是各 $X_j$ 所表示的\Term{functor}{函子}的\Term{projective-limit}{投射极限}，对这些\Term{functor}{函子}组成的\Term{projective-system}{投射系统}应用 9.6 (ii) 即可。
\end{proof}

\begin{SGAresult}{注记}{9.10}
在陈述 9.6、9.7、9.8 和 9.9 中，可以把所有“$\Pi$-\Term{accessible}{可达}”“\Term{accessible}{可达}”分别换成“$\Pi$-\Term{preaccessible}{预可达}”“\Term{preaccessible}{预可达}”。上述证明事实上也同时证明了这些陈述的这一变体。
\end{SGAresult}

\begin{SGAresult}{命题}{9.11}
设 $E$ 是一个满足条件 $\mathrm L$（9.1）的\Term{category}{范畴}，$C$ 是一个\Term{generating-subcategory}{小满生成子范畴}（7.1）。假定满足下述条件：

a) $E$ 对平行双\Term{arrow}{箭头}的\Term{equalizer}{等化子}\Term{stable}{稳定}，且 $E$ 的平行双\Term{arrow}{箭头}\Term{category}{范畴}上的\Term{functor}{函子} $\operatorname{Ker}$ 是\Term{accessible}{可达}的（例如，它与小\Term{filtered-inductive-limit}{滤过归纳极限}交换）。

则 $E$ 的每个\Term{object}{对象}都\Term{preaccessible}{预可达}，从而 $E$ 也\Term{preaccessible}{预可达}。进一步假定 $C$ \Term{generating-by-strict-epimorphisms}{由严格满态射生成}（7.1），且 $E$ 满足条件：

b) $E$ 对\Term{fiber-product}{纤维积}\Term{stable}{稳定}。
\SourcePage{91}
c) $E$ 中的每个小\Term{strict-epimorphic-family}{严格满态射族}都是\Term{universally-strict-epimorphic-family}{普适严格满态射族}（10.3）；或者，$E$ 中的小\Term{coproduct}{直和}\Term{representable}{可表示}，且 $E$ 的每个\Term{strict-epimorphism}{严格满态射}都是\Term{universally-strict-epimorphism}{普适严格满态射}。

则 $E$ 的每个\Term{object}{对象}都\Term{accessible}{可达}，从而 $E$ 也\Term{accessible}{可达}。
\end{SGAresult}

在条件 (a) 下 $E$ 的每个\Term{object}{对象}都\Term{accessible}{可达}，这来自以下事实：对于 $E$ 的每个\Term{object}{对象} $X$，$E$ 的\Term{strict-subobject}{严格子对象}集是小的（7.4），以及下述引理。

\begin{SGAresult}{引理}{9.11.1}
在 (9.11 a) 的条件下，设 $X\in\operatorname{ob}E$，$\Pi$ 是一个\Term{cardinal}{基数}，使得 $E$ 满足 $\mathrm L_\Pi$（9.1），(9.11 a) 中的\Term{functor}{函子} $\operatorname{Ker}$ 是 $\Pi$-\Term{accessible}{可达}的，且 $X$ 的\Term{strict-subobject}{严格子对象}集的\Term{cardinal}{基数}以 $\Pi$ 为上界。则 $X$ 是 $\Pi$-\Term{preaccessible}{预可达}的。
\end{SGAresult}
\begin{proof}
事实上，设 $I$ 是一个相对于 $I$ \Term{large-relative-to-cardinal}{足够大} 的\Term{set}{集合}，$(Y_i)_{i\in I}$ 是 $E$ 中的一个\Term{inductive-system}{归纳系统}，其\Term{inductive-limit}{归纳极限}为 $Y$；又设 $(u_i,v_i:X\rightrightarrows Y_i)$ 是一个平行双\Term{arrow}{箭头}，使得复合的平行双\Term{arrow}{箭头} $u,v:X\rightrightarrows Y$ 满足 $u=v$。证明 $I$ 中存在 $j\geq i$，使得 $u_j=v_j$。为此，考虑平行双\Term{arrow}{箭头}的\Term{inductive-system}{归纳系统} $(u_j,v_j:X\rightrightarrows Y_j)_{j\geq i}$，其\Term{inductive-limit}{归纳极限}为 $(u,v)$。由假设有
\[
X=\operatorname{Ker}(u,v)=\varinjlim_j\operatorname{Ker}(u_j,v_j)
=\varinjlim_jX_j,
\qquad X_j=\operatorname{Ker}(u_j,v_j).
\]
各 $X_j$ 都是 $X$ 的\Term{strict-subobject}{严格子对象}，故存在 $I$ 的一个\Term{subset}{子集} $I'$，由满足 $j\geq i$ 的指标组成，使得 $\operatorname{card}J\leq\Pi$，且每个 $X_j$ 都等于某个 $X_{i'}$（$i'\in I'$）。因为 $I$ 相对于 $\Pi$ \Term{large-relative-to-cardinal}{足够大}，$I'$ 在 $I$ 中有一个上界 $j$。于是 $X_j$ 包含所有 $j'\geq i$ \Term{correspondence}{对应}的 $X_{j'}$，故 $X_j\to X$ 是\Term{epimorphism}{满态射}（因为族 $X_{j'}\to X$ 是\Term{epimorphic-family}{满态射族}），又因它是\Term{strict-monomorphism}{严格单态射}，所以是\Term{isomorphism}{同构}。因而 $u_j=v_j$，这就证明了 9.11.1。
\end{proof}

9.11 的第二个断言来自第一个断言以及下述引理。

\begin{SGAresult}{引理}{9.11.2}
在 (9.11 a)、b)、c) 的条件下，设 $X$ 是 $E$ 的一个\Term{object}{对象}，$\Pi$ 是一个无限\Term{cardinal}{基数}，使得 $E$ 满足 $\mathrm L_\Pi$（9.1），$\operatorname{card}\operatorname{ob}C_{/X}\leq\Pi$，对 $C_{/X}$ 的两个\Term{object}{对象} $X'\to X$ 和 $X''\to X$，$X'\times_X X''$ 是 $\Pi$-\Term{preaccessible}{预可达}的，最后 $X$ 是 $\Pi$-\Term{preadmissible}{预可容许}的。则 $X$ 是 $\Pi$-\Term{accessible}{可达}的。
\end{SGAresult}
\begin{proof}
沿用 9.11.1 证明中的记号，只须证明每个\Term{morphism}{态射} $u:X\to Y$ 都来自一个\Term{morphism}{态射} $u_i:X\to Y_i$。考虑\Term{morphism}{态射}族 $Y_i\to Y$，它是\Term{strict-epimorphic-family}{严格满态射族}；由 b) 和 c)，族 $Y_i\times_YX\to X$ 也是\Term{strict-epimorphic-family}{严格满态射族}。另一方面，对于每个 $i$，因为 $C$ \Term{generating-by-strict-epimorphisms}{由严格满态射生成}，所以以 $C$ 中的\Term{object}{对象}为源的\Term{arrow}{箭头}
\[
T\longrightarrow Y_i\times_YX
\]
组成\Term{strict-epimorphic-family}{严格满态射族}。在 c) 所考虑的第二种情形下，可以找到这样的一个\Term{strict-epimorphism}{严格满态射}，其源是 $C$ 的\Term{object}{对象}的一个（小）和。由\Term{universally-strict-epimorphic-family}{普适严格满态射族}这一概念的\Term{transitivity}{传递性}（第 II 讲 2.5），
\SourcePage{92}
可知以 $C$ 中的\Term{object}{对象}为源、且经过某个 $Y_i\times_YX$ \Term{resolution}{分解}的\Term{arrow}{箭头} $f_\alpha:T_\alpha\to X$ 组成\Term{strict-epimorphic-family}{严格满态射族}。设 $J$ 是这个族的指标集，其\Term{cardinal}{基数}以 $\operatorname{card}\operatorname{ob}C_{/X}\leq\Pi$ 为上界。对每个 $\alpha\in J$，选取 $i=i(\alpha)\in I$ 和一个 $X$-\Term{morphism}{态射} $T_\alpha\to Y_i\times_YX$；\Term{equivalence}{等价}地，选取一个 $v_\alpha:T_\alpha\to Y_i$，使 $p_iv_\alpha=uf_\alpha$，其中 $p_i:Y_i\to Y$ 是\Term{canonical}{典范}\Term{morphism}{态射}。由于 $I$ 相对于 $\Pi$ \Term{large-relative-to-cardinal}{足够大}，可以使 $i(\alpha)$ 的选取与 $\alpha$ 无关，记为 $i$。对于每对指标 $\alpha,\beta\in J$，考虑复合
\[
\operatorname{pr}_1v_\alpha,\operatorname{pr}_2v_\beta:
T_\alpha\times_X T_\beta\rightrightarrows Y_i.
\]
它们与 $Y_i\to Y$ 的复合相等，所以，由于按假设 $T_\alpha\times_XT_\beta$ 是 $\Pi$-\Term{preaccessible}{预可达}的，存在指标 $i'=i(\alpha,\beta)\geq i$，使所考虑的\Term{arrow}{箭头}与 $Y_i\to Y_{i'}$ 的复合相等。所有对 $(\alpha,\beta)$ 组成的\Term{set}{集合}的\Term{cardinal}{基数} $\leq\Pi^2=\Pi$（因为 $\Pi$ 无限），故又可使 $i'$ 的选取与 $\alpha,\beta$ 无关。显然可以假定 $i'=i$。于是，因为族 $(f_\alpha:T_\alpha\to X)$ 是\Term{strict-epimorphic-family}{严格满态射族}，可以找到\Term{morphism}{态射} $u_i:X\to Y_i$，使得 $u_if_\alpha=v_i$。于是 $p_iu_i=p$，因为对每个 $\alpha$ 有
\[
(p_iu_i)f_\alpha=p_i(u_if_\alpha)=p_iv_\alpha=uf_\alpha,
\]
而族 $(f_\alpha)$ 是\Term{epimorphic-family}{满态射族}。这就完成了 9.11.2 的证明。
\end{proof}

\begin{SGAresult}{注记}{9.11.3}
作为 9.11 的特例，可知在以下两种情形中，\Term{category}{范畴} $E$ 的每个\Term{object}{对象}都\Term{accessible}{可达}：1) $E$ 是一个\Term{filtered-inductive-limit}{滤过归纳极限}\Term{exactness}{正合}、且具有小\Term{generating-subcategory}{生成子范畴}的 $\mathcal U$-\Term{abelian-category}{阿贝尔范畴}（这里适用 c) 的第二种情形）。2) $E$ 是\Term{topological-space}{拓扑空间} $X\in\mathcal U$ 上\Term{set}{集合}\Term{sheaf}{层}的\Term{category}{范畴}。更一般地，$E$ 是一个 $\mathcal U$-\Term{site}{景}上的\Term{set}{集合}\Term{sheaf}{层}\Term{category}{范畴}（第 II 讲 2.1）就够了，或者 $E$ 是一个 $\mathcal U$-\Term{topos}{拓扑斯}（第 IV 讲 1.1）也够了。
\end{SGAresult}

\begin{SGAresult}{定义}{9.12}
设 $E$ 是一个\Term{category}{范畴}。$E$ 的\Term{cardinal-filtration}{基数滤过}是指 $E$ 的一个由\Term{strictly-full-subcategory}{严格满子范畴} $\operatorname{Filt}^\Pi(E)$ 组成的递增\Term{filtered}{滤过} $(\operatorname{Filt}^\Pi(E))_{\Pi\geq\Pi_\circ}$，其指标为满足 $\Pi\geq\Pi_\circ$ 的\Term{cardinal}{基数} $\Pi\in\mathcal U$（其中 $\Pi_\circ$ 是一个固定的无限\Term{cardinal}{基数}，依赖于所考虑的\Term{cardinal-filtration}{基数滤过}），并满足下述条件：

a) 对每个 $\Pi\geq\Pi_\circ$，$\operatorname{Filt}^\Pi(E)$ 都\Term{equivalence}{等价}于一个\Term{small-category}{小范畴}。

b) $E$ 满足 $\mathrm L_{\Pi_\circ}$（9.3），且对每个 $\Pi\geq\Pi_\circ$，$\operatorname{Filt}^\Pi(E)$ 在 $E$ 中对以相对于 $\Pi_\circ$ \Term{large-relative-to-cardinal}{足够大}、且满足 $\operatorname{card}I\leq\Pi$ 的\Term{ordered-set}{有序集} $I$ 为指标的\Term{filtered-inductive-limit}{滤过归纳极限}\Term{stable}{稳定}。

c) 对每个 $\Pi\geq\Pi_\circ$ 和每个 $X\in\operatorname{ob}E$，可以找到一个\Term{isomorphism}{同构}
\[
X\simeq\varinjlim_I X_i,
\]
其中 $(X_i)_{i\in I}$ 是 $E$ 中以一个相对于 $\Pi$ \Term{large-relative-to-cardinal}{足够大} 的\Term{ordered-set}{有序集} $I$（9.1）为指标的\Term{inductive-system}{归纳系统}，所有 $X_i$ 都在 $\operatorname{Filt}^\Pi(E)$ 中。此外，若 $X\in\operatorname{ob}\operatorname{Filt}^\Pi(E)$，$\Pi'\geq\Pi$，则可以使 $I$ 满足 $\operatorname{card}I\leq\Pi'^\Pi$。
\end{SGAresult}

\SourcePage{93}
\textbf{9.12.1。}注意，由 b) 和 c) 可知，每个 $X\in\operatorname{ob}E$ 都属于某个 $\operatorname{Filt}^\Pi(E)$，其中 $\Pi$ 充分大。

\begin{SGAresult}{例}{9.12.2}
取 $E=(\mathcal U\text{-}\mathrm{Ens})$，$\Phi_\circ=\aleph_\circ$，$\operatorname{Filt}^\Pi(E)$ 为 $E$ 中由满足 $\operatorname{card}X\leq\Pi$ 的\Term{set}{集合}组成的\Term{full-subcategory}{满子范畴}。更一般地，有下述命题。
\end{SGAresult}

\begin{SGAresult}{命题}{9.13}
设 $E$ 是一个 $\mathcal U$-\Term{category}{范畴}。假定 $E$ 对小\Term{filtered-inductive-limit}{滤过归纳极限}、两个\Term{object}{对象}的和及平行双\Term{arrow}{箭头}的\Term{coequalizer}{余等化子}\Term{stable}{稳定}，$E$ 对\Term{fiber-product}{纤维积}\Term{stable}{稳定}，且 $E$ 中的\Term{epimorphism}{满态射}都是\Term{universally-strict-epimorphism}{普适严格满态射}。设 $C$ 是 $E$ 的一个\Term{generating-by-strict-epimorphisms}{由严格满态射生成}的小\Term{full-subcategory}{满子范畴}。设 $\Pi_\circ$ 是一个满足 $\Pi_\circ\geq\operatorname{card}\operatorname{Fl}C$ 的无限\Term{cardinal}{基数}。对每个\Term{cardinal}{基数} $\Pi\geq\Pi_\circ$，以 $\operatorname{Filt}^\Pi(E)$ 表示 $E$ 中由下述\Term{object}{对象} $X$ 组成的\Term{strictly-full-subcategory}{严格满子范畴}：存在一个以 $X$ 为目标的\Term{strict-epimorphic-family}{严格满态射族} $(X_i\to X)_{i\in I}$，使得 $\operatorname{card}I\leq\Pi$，且对每个 $i\in I$ 都有 $X_i\in\operatorname{ob}C$。则 $(\operatorname{Filt}^\Pi(E))_{\Pi\geq\Pi_\circ}$ 是 $E$ 的一个\Term{cardinal-filtration}{基数滤过}。此外，对每个 $X\in\operatorname{ob}\operatorname{Filt}^\Pi(E)$，$C_{/X}$ 的\Term{arrow}{箭头}集（从而其\Term{object}{对象}集）的\Term{cardinal}{基数}以 $\Pi^{\Pi_\circ}$ 为上界。
\end{SGAresult}
\begin{proof}
最后一个断言就是 7.6。

为说明我们得到一个\Term{cardinal-filtration}{基数滤过}，必须证明 9.12 的条件 a)、b)、c)。条件 a) 立即由 7.5.2 得到。对于 b)，设 $X=\varinjlim_I X_i$，其中 $\operatorname{card}I\leq\Pi$，且各 $X_i\in\operatorname{ob}\operatorname{Filt}^\Pi(E)$。于是\Term{canonical}{典范}\Term{morphism}{态射}族 $X_i\to X$ 是\Term{strict-epimorphic-family}{严格满态射族}，且由关于 $X_i$ 的假设，对每个 $i\in I$，存在一个\Term{strict-epimorphic-family}{严格满态射族} $(X_{ij}\to X_i)_{j\in J_i}$，其中 $\operatorname{card}J_i\leq\Pi$。转而考虑和 $\coprod_iX_i$ 与 $\coprod_{i,j}X_{ij}$，由\Term{universally-strict-epimorphism}{普适严格满态射}的\Term{transitivity}{传递性}（第 II 讲 2.5）可知，复合 $X_{ij}\to X_i\to X$ 所成的族（以所有 $i\in I$ \Term{correspondence}{对应}的 $J_i$ 的\Term{sum-set}{和集}为指标）是\Term{strict-epimorphic-family}{严格满态射族}。而 $\operatorname{card}J\leq\Pi^2=\Pi$，故结论成立。（注意：这里只用了 $\varinjlim_I X_i$ 存在，并没有用 $I$ 相对于 $\Pi_\circ$ \Term{large-relative-to-cardinal}{足够大}，甚至没有用它是\Term{filtered-category}{滤过的}。）

最后证明 c)。注意，对每个 $X\in\operatorname{ob}E$，以 $\operatorname{ob}C$ 中的\Term{object}{对象}为源的\Term{arrow}{箭头} $X_i\to X$ 所成的族是\Term{strict-epimorphic-family}{严格满态射族}，因为 $C$ \Term{generating-by-strict-epimorphisms}{由严格满态射生成}；这个族显然是小的，故存在\Term{cardinal}{基数} $\Pi\geq\Pi_\circ$，使 $X\in\operatorname{ob}\operatorname{Filt}^\Pi(E)$。剩下只须证明：若有\Term{cardinal}{基数} $\Pi'\geq\Pi\geq\Pi_\circ$，且 $X\in\operatorname{ob}\operatorname{Filt}^{\Pi'}(E)$，则有\Term{isomorphism}{同构}
\[
X\simeq\varinjlim_I X_i,
\]
其中 $I$ 相对于 $\Pi$ \Term{large-relative-to-cardinal}{足够大}，$\operatorname{card}I\leq\Pi'^\Pi$，且各 $X_i$ 在 $\operatorname{ob}\operatorname{Filt}^\Pi(E)$ 中。因为 $C$ \Term{generating-by-strict-epimorphisms}{由严格满态射生成}，所以
\[
X\simeq\varinjlim_{C_{/X}}T.\tag{*}
\]
还须注意，依次向 $C$ 加入 $C$ 的\Term{object}{对象}的和及平行双\Term{arrow}{箭头}的\Term{coequalizer}{余等化子}，再对所得\Term{category}{范畴}作同样操作，依此类推，便归结到 $C$ 在 $E$ 中对两个\Term{object}{对象}的和及平行双\Term{arrow}{箭头}的\Term{coequalizer}{余等化子}\Term{stable}{稳定}的情形，
\SourcePage{94}
而且这不会破坏假设 $\Pi_\circ\geq\operatorname{card}\operatorname{Fl}C$。还可假定，若 $E$ 含有一个选定的\Term{initial-object}{始对象} $\varnothing_E$，就有 $\varnothing_E\in\operatorname{ob}C$。这蕴涵\Term{category}{范畴} $C_{/X}$ 对两个\Term{object}{对象}的和及平行双\Term{arrow}{箭头}的\Term{coequalizer}{余等化子}\Term{stable}{稳定}。于是它是\Term{filtered-category}{滤过的}，因为它\Term{nonempty}{非空}：若它为空，关系 (*) 将表明 $X$ 是\Term{initial-object}{始对象}，因而 $C_{/X}$ 含有\Term{arrow}{箭头} $\varnothing_E\to X$，矛盾。

令 $I$ 为 $C_{/X}$ 的\Term{full-subcategory}{满子范畴} $i$ 组成的\Term{set}{集合}，其中 $i$ 是\Term{filtered-category}{滤过的}，且 $\operatorname{card}\operatorname{ob}i\leq\Pi$，并按包含关系排序。对每个 $i\in I$，令 $X_i\in\operatorname{ob}E$ 为复合\Term{functor}{函子} $i\to C_{/X}\to E$ 的\Term{inductive-limit}{归纳极限}，它由假设存在。显然有
\[
\varinjlim_I X_i\simeq\varinjlim_{C_{/X}}T\simeq X.
\]
另一方面，已经指出 $\operatorname{card}\operatorname{ob}C_{/X}\leq\Pi'^{\Pi_\circ}$。由下述引理可知 $I$ 相对于 $\Pi$ \Term{large-relative-to-cardinal}{足够大}，且
\[
\operatorname{card}I\leq(\Pi'^{\Pi_\circ})^\Pi
=\Pi'^{\Pi\Pi_\circ}=\Pi'^\Pi.
\]
该引理的证明留给读者（其中取 $I=C_{/X}$，$c=\Pi'^{\Pi_\circ}$）。
\end{proof}

\begin{SGAresult}{引理}{9.13.1}
设 $J$ 是一个\Term{filtered}{滤过}\Term{category}{范畴}，$c,\Pi$ 是两个\Term{cardinal}{基数}，满足 $c\geq\sup(\operatorname{card}\operatorname{Fl}J,\Pi)$，且对 $J$ 的每对\Term{object}{对象} $j,j'$，有 $\operatorname{card}\operatorname{Hom}(j,J')\leq\Pi_\circ$。设 $I$ 是 $J$ 的所有满足 $\operatorname{card}\operatorname{ob}i\leq\Pi$ 的\Term{filtered-category}{滤过满子范畴} $i$ 所成的\Term{set}{集合}。则按包含关系排序后，$I$ 相对于 $\Pi$ \Term{large-relative-to-cardinal}{足够大}，且 $\operatorname{card}I\leq c^\Pi$。
\end{SGAresult}

这就完成了 9.13 的证明。

\begin{SGAresult}{注记}{9.13.2}
略作进一步努力，应当可以在 9.13 中将“$E$ 对小\Term{filtered-inductive-limit}{滤过归纳极限}\Term{stable}{稳定}”的假设换成“$E$ 满足 $\mathrm L$”（9.1），只要假定 $E$ 对小和\Term{stable}{稳定}，或者 $E$ 中的每个\Term{epimorphic-family}{满态射族}都是\Term{universally-epimorphic-family}{普适满态射族}。在 c) 的证明中，此时应选取 $\Pi_\circ$ 使 $E$ 满足 $\mathrm L_{\Pi_\circ}$，并且只考虑 $C_{/X}$ 中不仅是\Term{filtered-category}{滤过的}、而且其\Term{preordered-set}{预序集} $\operatorname{ob}i$ 相对于 $\Pi_\circ$ \Term{large-relative-to-cardinal}{足够大} 的\Term{full-subcategory}{满子范畴} $i$。利用第 8 节以及 $E$ 满足 $\mathrm L_{\Pi_\circ}$ 的假设，可知各 $X_i$ 存在，然后应可借助 9.13.1 的一个适当变体得出结论，不过执笔者尚未核实这个变体。
\end{SGAresult}

\begin{SGAresult}{命题}{9.14}
设 $f:E\to F$ 是两个带有\Term{cardinal-filtration}{基数滤过} $(\operatorname{Filt}^\Pi(E))_{\Pi\geq\Pi_\circ}$ 和 $(\operatorname{Filt}^\Pi(F))_{\Pi\geq\Pi'_\circ}$ 的\Term{category}{范畴}之间的一个\Term{accessible-functor}{可达函子}（9.2）。则存在\Term{cardinal}{基数} $\Pi_1\geq\sup(\Pi_\circ,\Pi'_\circ)$，使对每个\Term{cardinal}{基数} $\Pi\geq\Pi_1$，有
\[
f(\operatorname{Filt}^\Pi(E))\subset\operatorname{Filt}^{\Pi^{\Pi'_\circ}}(F).\tag{9.14.1}
\]
特别地，若 $\Pi=2^c$，其中 $c\geq\Pi_1$，从而 $\Pi^{\Pi_1}=2^{c\Pi_1}=2^c=\Pi$，则有
\[
f(\operatorname{Filt}^\Pi(E))\subset\operatorname{Filt}^\Pi(F).\tag{9.14.2}
\]
\end{SGAresult}

立即指出下述推论。

\SourcePage{95}
\begin{SGAresult}{推论}{9.15}
设 $E$ 是一个带有两个\Term{cardinal-filtration}{基数滤过} $(\operatorname{Filt}^\Pi(E))_{\Pi\geq\Pi_\circ}$ 与 $(\operatorname{Filt}'^\Pi(E))_{\Pi\geq\Pi'_\circ}$ 的\Term{category}{范畴}。则存在\Term{cardinal}{基数} $\Pi_1\geq\sup(\Pi_\circ,\Pi'_\circ)$，使得对每个\Term{cardinal}{基数} $c\geq\Pi_1$，置 $\Pi=2^c$，便有 $\operatorname{Filt}^\Pi(E)=\operatorname{Filt}'^\Pi(E)$。
\end{SGAresult}

\begin{proof}[9.14 的证明]
设 $c$ 是一个满足 $c\geq\sup(\Pi_\circ,\Pi'_\circ)$ 的\Term{cardinal}{基数}，并使 $f$ 是 $c$-\Term{admissible}{可容许}的。又设 $\Pi_1\geq c$ 使得
\[
f(\operatorname{Filt}^c(E))\subset\operatorname{Filt}^{\Pi_1}(F).\tag{*}
\]
这样的 $\Pi_1$ 存在，因为 $\operatorname{Filt}^c(E)$ \Term{equivalence}{等价}于一个\Term{small-category}{小范畴}（9.12 a)），所以 $f(\operatorname{Filt}^c(E))$ 也是如此。因而可以对后者的\Term{object}{对象}应用 9.12.1，找到一个包含它们全体的 $\operatorname{Filt}^\Pi(F)$（这里用到各 $\operatorname{Filt}^\Pi(F)$ 都是\Term{full-subcategory}{满子范畴}）。于是设 $\Pi\geq\Pi_1$，$X\in\operatorname{ob}\operatorname{Filt}^\Pi(E)$，证明 $f(X)\in\operatorname{Filt}^{\Pi^{\Pi_1}}(F)$。事实上，写成
\[
X=\varinjlim_I X_i,
\]
其中各 $X_i\in\operatorname{ob}\operatorname{Filt}^c(E)$，$I$ 相对于 $c$ \Term{large-relative-to-cardinal}{足够大}，且 $\operatorname{card}I\leq\Pi^c\leq\Pi^{\Pi_1}$（9.12 c)）。由于 $f$ 是 $c$-\Term{admissible}{可容许}的，$I$ 相对于 $c$ \Term{large-relative-to-cardinal}{足够大}，有
\[
f(X)\simeq\varinjlim_I f(X_i),
\]
又因 $\operatorname{card}I\leq\Pi^{\Pi_1}$，且由 (*) 有
\[
f(X_i)\in\operatorname{ob}\operatorname{Filt}^{\Pi_1}(F)
\subset\operatorname{ob}\operatorname{Filt}^{\Pi^{\Pi_1}}(F),
\]
所以由 9.12 b)，$f(X)\in\operatorname{Filt}^{\Pi^{\Pi_1}}(F)$，证毕。
\end{proof}

\textbf{9.15.1。}\Term{cardinal-filtration}{基数滤过}（9.12）这一概念，主要在 $E$ 的\Term{object}{对象}都\Term{accessible}{可达}时才有意义。指出，由 9.11 可知，如果在 9.13 的假设之外再假定 $E$ 的平行双\Term{arrow}{箭头}\Term{category}{范畴}上的\Term{functor}{函子} $\operatorname{Ker}$ \Term{accessible}{可达}，则这一条件成立。另一方面还有：

\begin{SGAresult}{命题}{9.16}
设 $E$ 是一个带有\Term{cardinal-filtration}{基数滤过} $(\operatorname{Filt}^\Pi(E))_{\Pi\geq\Pi_\circ}$ 的\Term{category}{范畴}。假定 $\operatorname{Filt}^{\Pi_\circ}(E)$ 的元素都是 $E$ 的\Term{accessible-object}{可达对象}；则存在 $\mathcal U$ 中的一个\Term{cardinal}{基数} $\Pi_1\geq\Pi_\circ$，使 $\operatorname{Filt}^{\Pi_\circ}(E)$ 的\Term{object}{对象}都是 $\Pi_1$-\Term{accessible}{可达}的。如此选取 $\Pi_1$ 后，对每个\Term{cardinal}{基数} $\Pi\geq\Pi_1$，有（沿用 9.3.1 的记号）
\[
E_\Pi\subset\operatorname{Filt}^\Pi(E)\subset E_{(\Pi^{\Pi_\circ})}.\tag{9.16.1}
\]
\end{SGAresult}
\begin{proof}
$\Pi_1$ 的存在性直接来自 $\operatorname{Filt}^{\Pi_\circ}(E)$ \Term{equivalence}{等价}于一个\Term{small-category}{小范畴}（9.12 a)）。设 $\Pi\geq\Pi_1$。若 $X$ 在 $\operatorname{Filt}^\Pi(E)$ 中，将它写成 $X\simeq\varinjlim_I X_i$，其中 $\operatorname{card}I\leq\Pi^{\Pi_\circ}$，对每个 $i$ 都有 $X_i\in\operatorname{ob}\operatorname{Filt}^{\Pi_\circ}(E)$（9.12 c)），则由 9.9 可知 $X$ 是 $\Pi^{\Pi_\circ}$-\Term{accessible}{可达}的，故得 (9.16.1) 的第二个包含关系。

假定 $X$ 是 $\Pi$-\Term{accessible}{可达}的，写成 $X\simeq\varinjlim_I X_i$，其中 $I$ 相对于 $\Pi$ \Term{large-relative-to-cardinal}{足够大}，且各 $X_i\in\operatorname{ob}\operatorname{Filt}^\Pi(E)$（9.12 c)）。由关于 $X$ 的假设，给定的\Term{isomorphism}{同构} $X\xrightarrow{\sim}\varinjlim_I X_i$ 经过某个 $X_i$ \Term{resolution}{分解}，
\SourcePage{96}
所以 $X$ \Term{isomorphism}{同构}于这个 $X_i$ 的一个\Term{direct-factor}{直因子}。由下述引理便知 $X\in\operatorname{ob}\operatorname{Filt}^\Pi(E)$，从而得到 (9.16.1) 的第一个包含关系。
\end{proof}

\begin{SGAresult}{引理}{9.16.2}
$E$ 中任意一个是 $\operatorname{Filt}^\Pi(E)$ 的某个\Term{object}{对象}的\Term{direct-factor}{直因子}的\Term{object}{对象}，都在 $\operatorname{Filt}^\Pi(E)$ 中。
\end{SGAresult}
\begin{proof}
事实上，若 $X$ 是 $E$ 的\Term{object}{对象} $Y$ 上一个\Term{projector}{投影子} $p$ 的像（即 $p$ 是满足 $p^2=p$ 的\Term{endomorphism}{自同态}），且 $I$ 是一个\Term{filtered-category}{滤过有序集}，则 $X$ 是\Term{filtered}{滤过}\Term{inductive-system}{归纳系统} $(Y_i)_{i\in I}$ 的\Term{inductive-limit}{归纳极限}，其中对每个 $i\in I$ 令 $Y_i=Y$，当 $i<j$ 时令 $p:Y_i\to Y_j$。取 $I$ 相对于 $\Pi_\circ$ \Term{large-relative-to-cardinal}{足够大}，且 $\operatorname{card}I\leq\Pi$，便由 9.12 b) 看出，若 $Y\in\operatorname{Filt}^\Pi(E)$，则 $X$ 也是如此。
\end{proof}

\begin{SGAresult}{推论}{9.17}
在 9.16 的条件下，对每个\Term{cardinal}{基数} $c\geq\Pi_1$，置 $\Pi=2^c$，则 $\operatorname{Filt}^\Pi(E)$ 就是 $E$ 中由 $\Pi$-\Term{accessible-object}{可达对象}组成的\Term{strictly-full-subcategory}{严格满子范畴} $E_\Pi$。
\end{SGAresult}

\begin{SGAresult}{推论}{9.18}
设 $E$ 是一个满足条件 $\mathrm L_\Pi$（9.1）的\Term{category}{范畴}，其中 $\Pi$ 是一个\Term{cardinal}{基数}，$C$ 是 $E$ 的一个\Term{full-subcategory}{满子范畴}。以 $\operatorname{Ind}(C)_\Pi$ 表示 $C$ 的 \Term{ind-object}{ind-对象}\Term{category}{范畴} $\operatorname{Ind}(C)$（8.2）中由形如 $(X_i)_{i\in I}$ 的 \Term{ind-object}{ind-对象}组成的\Term{full-subcategory}{满子范畴}，其中 $I$ 是一个相对于 $\Pi$ \Term{large-relative-to-cardinal}{足够大} 的\Term{ordered-set}{有序集}。考虑\Term{canonical}{典范}\Term{functor}{函子}
\[
(X_i)_{i\in I}\longmapsto\varinjlim_I X_i:
\operatorname{Ind}(C)_\Pi\longrightarrow E.\tag{9.18.1}
\]
a) 这个\Term{functor}{函子}\Term{fully-faithful-functor}{全忠实}的充要条件是 $C$ 的每个\Term{object}{对象} $X$ 都是 $E$ 的 $\Pi$-\Term{accessible-object}{可达对象}（9.3）。

b) 在 9.17 的条件下，特别是 $\Pi=2^c$，取 $C=\operatorname{Filt}^\Pi(E)$。则\Term{functor}{函子} (9.18.1) 是\Term{equivalence-of-categories}{范畴等价}。
\end{SGAresult}
\begin{proof}
断言 a) 是 8.7.5 a) 的直接推广，证明方式也相同。于是 b) 由 9.17 以及\Term{cardinal-filtration}{基数滤过}的条件 9.12 c) 得到，后者蕴涵所考虑的\Term{functor}{函子}是\Term{essentially-surjective-functor}{本质满射}的。
\end{proof}

\begin{SGAresult}{推论}{9.19}
在 9.17 的条件下，设 $C=\operatorname{Filt}^\Pi(E)$，$F$ 是一个 $\mathcal U$-\Term{category}{范畴}，考虑由“\Term{restriction}{限制}到 $C$”\Term{functor}{函子} $f\mapsto f|_C$ 诱导的\Term{functor}{函子}
\[
\mathscr{H}om(E,F)_\Pi\longrightarrow\mathscr{H}om(C,F),\tag{9.19.1}
\]
其中 (9.19.1) 的源是从 $E$ 到 $F$ 的 $\Pi$-\Term{accessible-functor}{可达函子}的\Term{category}{范畴}（9.2）。上述\Term{functor}{函子}是\Term{fully-faithful-functor}{全忠实}的。若 $F$ 满足条件 $\mathrm L_\Pi$（9.1），则\Term{functor}{函子} (9.19.1) 是\Term{equivalence-of-categories}{范畴等价}。
\end{SGAresult}
\begin{proof}
第一个断言的证明与 7.8 相同，利用 9.18 b)。第二个断言通过构造 (9.19.1) 的\Term{quasi-inverse-functor}{拟逆函子}得到：对每个\Term{functor}{函子} $f:C\to F$，关联从 $\operatorname{Ind}(E_\Pi)$ 到 $F$ 的\Term{functor}{函子} $(X_i)_{i\in I}\mapsto\varinjlim_i f(X_i)$，
\SourcePage{97}
再利用 9.18 b) 从中得到\Term{functor}{函子} $f:E\to F$。一切归结为说明后者是 $\Pi$-\Term{accessible}{可达}的；证明与类似断言 8.7.3 的证明相同。
\end{proof}

\begin{SGAresult}{推论}{9.20}
设 $E$ 是一个具有\Term{cardinal-filtration}{基数滤过}、且每个\Term{object}{对象}都\Term{accessible}{可达}的\Term{category}{范畴}（参见 9.16），$F$ 是一个\Term{category}{范畴}。则：

a) 从 $E$ 到 $F$ 的\Term{accessible-functor}{可达函子}的\Term{category}{范畴} $\mathscr{H}om(E,F)_{\mathrm{acc}}$（9.2）是一个 $\mathcal U$-\Term{category}{范畴}。（注意，回顾 9.0，给定的\Term{category}{范畴} $E,F$ 都假定为 $\mathcal U$-\Term{category}{范畴}。）

b) 假定 $F$ 对小\Term{filtered-inductive-limit}{滤过归纳极限}\Term{stable}{稳定}。对每个\Term{equivalence}{等价}于一个\Term{small-category}{小范畴}的 $E$ 的\Term{full-subcategory}{满子范畴} $C$，设\Term{inclusion-functor}{包含函子}为 $i:C\to E$，考虑\Term{correspondence}{对应}的\Term{functor}{函子}
\[
i_!:\mathscr{H}om(C,F)\longrightarrow\mathscr{H}om(E,F)
\]
（5.1）。\Term{functor}{函子} $f:E\to F$ \Term{accessible}{可达}的充要条件是存在 $E$ 的一个小\Term{subcategory}{子范畴} $C$，使得 $f$ 属于上述\Term{functor}{函子} $i_!$ 的\Term{essential-image}{本质像}。
\end{SGAresult}
\begin{proof}
a) 只须证明，对每个使 $E$ 满足 $\mathrm L_\Pi$ 的\Term{cardinal}{基数} $\Pi$，$\mathscr{H}om(E,F)_{\mathrm{acc}}$ 的\Term{full-subcategory}{满子范畴} $\mathscr{H}om(E,F)_\Pi$ 是一个 $\mathcal U$-\Term{category}{范畴}。显然只须对形如 $2^c$、其中 $c$ 充分大的\Term{cardinal}{基数}核实这一点。但这由 9.19 得到，因为 $C=\operatorname{Filt}^\Pi(E)$ 是\Term{essentially-small-category}{本质小}的，$\mathscr{H}om(C,F)$ 显然是一个 $\mathcal U$-\Term{category}{范畴}。

b) 由\Term{functor}{函子} $i_!$ 的构成的\Term{transitivity}{传递性}，陈述中只须考虑形如 $\operatorname{Filt}^\Pi(E)$ 的\Term{subcategory}{子范畴} $C$，其中 $\Pi$ 如 9.17。此时容易看出，复合\Term{functor}{函子}
\[
\mathscr{H}om(\operatorname{Filt}^\Pi(E),F)
\longrightarrow\mathscr{H}om(E,F)_\Pi
\longrightarrow\mathscr{H}om(E,F),
\]
其中第一个\Term{arrow}{箭头}是 (9.19.1) 的\Term{quasi-inverse-functor}{拟逆}，第二个\Term{arrow}{箭头}是包含，恰好就是\Term{functor}{函子} $i_!$，相差一个\Term{isomorphism}{同构}。因此断言 b) 由 9.19 得到。
\end{proof}

\begin{SGAresult}{*习题}{9.20.1}
（本习题使用第 II 讲和第 IV 讲中发展的\Term{site}{景}与\Term{topos}{拓扑斯}的概念。）设 $E$ 是一个 $\mathcal U$-\Term{topos}{拓扑斯}，$\mathcal V$ 是一个满足 $\mathcal U\in\mathcal V$ 的\Term{universe}{宇宙}，$C$ 是 $\mathcal U$-\Term{topos}{拓扑斯} $E$ 的一个\Term{generating-subcategory}{小满生成子范畴}，$\Pi_\circ\in\mathcal U$ 是一个满足 $\Pi_\circ\in\operatorname{card}\operatorname{Fl}C$ 的无限\Term{cardinal}{基数}。对每个\Term{cardinal}{基数} $\Pi\geq\Pi_\circ$、$\Pi\in\mathcal U$，以 $\operatorname{Filt}^\Pi(E)$ 表示 $E$ 中由下述\Term{object}{对象} $X$ 组成的\Term{strictly-full-subcategory}{严格满子范畴}：存在一个以 $X$ 为目标的\Term{strict-epimorphic-family}{严格满态射族} $(X_i\to X)_{i\in I}$，使得 $\operatorname{card}I\leq\Pi$，且对每个 $i\in I$，$X_i\in\operatorname{Ob}C$。

a) 证明 $(\operatorname{Filt}^\Pi(E))_{\Pi\geq\Pi_\circ}$ 是 $\mathcal U$-\Term{category}{范畴} $E$ 的一个\Term{cardinal-filtration}{基数滤过}，而且可以选取 $\Pi_\circ$，使得对每个 $\Pi\geq\Pi_\circ$，有包含关系
\[
E_\Pi\subset\operatorname{Filt}^\Pi(E)\subset E_{\Pi^{\Pi_\circ}},
\]
其中对每个\Term{cardinal}{基数} $c$，$E_c$ 表示 $E$ 中由 $c$-\Term{accessible-object}{可达对象}组成的\Term{strictly-full-subcategory}{严格满子范畴}。

\SourcePage{98}
b) 选取一个满足 $\Pi=\Pi^{\Pi_\circ}$ 的\Term{cardinal}{基数} $\Pi\geq\Pi_\circ$（例如形如 $2^c$，其中 $c>\Pi_\circ$），并置 $C=\operatorname{Filt}^\Pi(E)$，证明有\Term{equivalence-of-categories}{范畴等价}
\[
\operatorname{Ind}(C)_\Pi\xrightarrow{\approx}E
\]
（$\operatorname{Ind}(C)_\Pi$ 的记号见 9.18）。

c) 保持 b) 的记号，设 $\mathcal V$ 是一个满足 $\mathcal U\subset\mathcal V$ 的\Term{universe}{宇宙}，$\widetilde E_{\mathcal V}$ 是\Term{site}{景} $E$（赋予其\Term{canonical-topology}{典范拓扑}）上的 $\mathcal V$-\Term{sheaf}{层}的 $\mathcal V$-\Term{topos}{拓扑斯}。以 $\operatorname{Ind}(C,\mathcal V)_\Pi$ 表示 $C$ 中以相对于 $\Pi$ \Term{large-relative-to-cardinal}{足够大} 的预序指标集为指标的 $\mathcal V$-\Term{ind-object}{ind-对象}的\Term{category}{范畴}。证明有\Term{equivalence-of-categories}{范畴等价}
\[
\operatorname{Ind}(C,\mathcal V)_\Pi\xrightarrow{\approx}\widetilde E_{\mathcal V}.
\]

d) 设 $c\in\mathcal V$ 是一个\Term{cardinal}{基数}，$\operatorname{Ind}(E,\mathcal V)'_c$ 是 $\operatorname{Ind}(E,\mathcal V)$ 中由以一个相对于每个\Term{cardinal}{基数} $<c$ 都\Term{large-relative-to-cardinal}{足够大} 的\Term{preordered-set}{预序集}为指标的 $E$ 的 $\mathcal V$-\Term{ind-object}{ind-对象}组成的\Term{full-subcategory}{满子范畴}。取 $c=\operatorname{card}\mathcal U$，证明有\Term{equivalence-of-categories}{范畴等价}
\[
\operatorname{Ind}(E,\mathcal V)'_c\xrightarrow{\approx}\widetilde E_{\mathcal V}.
\]
\hfill *
\end{SGAresult}

\textbf{9.21。}本款 9.21 为下述定理 9.22 的证明作技术准备，后者是本节第 9 节的主要结果。设
\[
p:E\longrightarrow B\tag{9.21.1}
\]
是一个\Term{fibred-functor}{纤维化函子}，其中 $B$ 是一个\Term{small-category}{小范畴}，且与 $B$ 的\Term{arrow}{箭头} $f:\alpha\to\beta$ 关联的\Term{inverse-image-functor}{逆像函子}
\[
f^*:E_\beta\longrightarrow E_\alpha
\]
是\Term{accessible}{可达}的（9.2）。特别地，各\Term{fiber-category}{纤维范畴} $E_\alpha$（$\alpha\in\operatorname{ob}B$）满足条件 $\mathrm L$（9.1）；所以，由于 $B$ 是小的，存在一个\Term{cardinal}{基数} $\Pi'_\circ\in\mathcal U$，使所有\Term{category}{范畴} $E_\alpha$ 都满足条件 $\mathrm L_{\Pi'_\circ}$。设 $\Pi_\circ\in\mathcal U$ 是一个无限\Term{cardinal}{基数}，且 $\Pi_\circ>\Pi'_\circ$，因而有
\[
\Pi_\circ>\Pi'_\circ,\qquad
E_\alpha\text{ 对每个 }\alpha\in\operatorname{ob}B\text{ 满足 }\mathrm L_{\Pi'_\circ}.
\tag{9.21.2}
\]
此外，所作假设立即蕴涵存在一个\Term{cardinal}{基数} $c\in\mathcal U$，满足下述条件：
\[
\left\{\begin{aligned}
{}&\text{a) }c\text{ 是无限的，}\\
&\text{b) }c\geq\operatorname{card}\operatorname{Fl}B,\\
&\text{c) 对 }B\text{ 中每个箭头 }f:\alpha\to\beta,
\ f^*:E_\beta\to E_\alpha\text{ 是 }c\text{-可达的。}
\end{aligned}\right.\tag{9.21.3}
\]
进一步假定，对每个 $\alpha\in\operatorname{ob}B$，可以找到 $E_\alpha$ 的一个\Term{full-subcategory}{满子范畴} $C_\alpha$，
\SourcePage{99}
满足下述条件：
\[
\left\{\begin{aligned}
{}&\text{a) 对每个 }X\in\operatorname{ob}C_\alpha,
\ X\text{ 在 }E_\alpha\text{ 中是 }c\text{-可达的（9.3）。}\\
&\text{b) 对每个 }X\in\operatorname{ob}E_\alpha,
\text{ 可在 }E_\alpha\text{ 中找到同构 }X\simeq\varinjlim_I X_i,\\
&\qquad\text{其中各 }X_i\text{ 在 }C_\alpha\text{ 中，且 }I\text{ 是一个大于 }c\text{ 的有序集。}
\end{aligned}\right.\tag{9.21.4}
\]
设 $\Pi$ 是满足下式的一个\Term{cardinal}{基数}：
\[
\Pi\geq\Pi_\circ,\qquad\text{从而 }\Pi>\Pi'_\circ,\tag{9.21.5}
\]
对每个 $\alpha\in\operatorname{ob}B$，令
\[
C_\alpha^\Pi\subset E_\alpha\tag{9.21.6}
\]
由可以表示成 $\varinjlim_I X_i$ 的\Term{object}{对象}组成，其中 $I$ 是一个相对于 $\Pi'_\circ$ \Term{large-relative-to-cardinal}{足够大}、满足 $\operatorname{card}I\leq\Pi$ 的\Term{ordered-set}{有序集}，且各 $X_i$ 都在 $C_\alpha$ 中。由 7.5.2 显然可知，$C_\alpha^\Pi$ 是\Term{essentially-small-category}{本质小}的，即\Term{equivalence}{等价}于一个\Term{small-category}{小范畴}。

设
\[
F=\mathscr{H}om_B(B,E)\tag{9.21.7}
\]
为 $E$ 在 $B$ 上的\Term{section-category}{截面范畴}，令
\[
F^\Pi\subset F\tag{9.21.8}
\]
为 $F$ 中由满足对每个 $\alpha\in\operatorname{ob}B$ 都有
\[
X(\alpha)\in C_\alpha^\Pi
\]
的\Term{section}{截面} $X:\alpha\mapsto X(\alpha)$ 组成的\Term{strictly-full-subcategory}{严格满子范畴}。

由于各 $C_\alpha^\Pi$ 是\Term{essentially-small-category}{本质小}的，\Term{category}{范畴} $F^\Pi$ 显然也是如此。我们将说明这个\Term{category}{范畴}是生成的，更精确地说：

\begin{SGAresult}{引理}{9.21.9}
在上述条件和记号下，有下述结论：

(i) $F$ 的每个\Term{object}{对象}都\Term{accessible}{可达}（9.3）。若 $d$ 是一个\Term{cardinal}{基数}，$d\geq\Pi'_\circ$，且 $\alpha\mapsto X(\alpha)$ 是 $F$ 的一个元素，使得对每个 $\alpha$，$X(\alpha)$ 在 $E_\alpha$ 中是 $d$-\Term{accessible}{可达}的，则 $X$ 在 $F$ 中是 $d$-\Term{accessible}{可达}的。

(ii) 假定 $\Pi\leq c$，或者 $\Pi'_\circ$ \Term{finite}{有限}（即各 $E_\alpha$ 对小\Term{filtered-inductive-limit}{滤过归纳极限}\Term{stable}{稳定}）。则 $F$ 的每个\Term{object}{对象} $X$ 都\Term{isomorphism}{同构}于形如 $\varinjlim_I X_i$ 的\Term{object}{对象}，其中各 $X_i$ 都在 $F^\Pi$ 中，且 $I$ 相对于 $\Pi$ \Term{large-relative-to-cardinal}{足够大}。
\end{SGAresult}

为证明 (i)，先注意由 (9.21.4) 和 9.9，$E_\alpha$ 的每个\Term{object}{对象}都\Term{accessible}{可达}。另一方面，由下述引理，$F$ 满足条件 $\mathrm L_{\Pi'_\circ}$。

\SourcePage{100}
\begin{SGAresult}{引理}{9.21.10}
设 $p:E\to B$ 是一个\Term{fibred-functor}{纤维化函子}，$F=\mathscr{H}om_B(B,E)$，$I$ 是一个\Term{category}{范畴}，$i\mapsto X_i$ 是从 $I$ 到 $F$ 的\Term{functor}{函子}。$\varinjlim_I X_i$ 在 $F$ 中\Term{representable}{可表示}的一个充分条件是，对每个 $\alpha\in\operatorname{ob}B$，$\varinjlim_I X_i(\alpha)$ 在\Term{fiber-category}{纤维范畴} $E_\alpha$ 中\Term{representable}{可表示}；此时，$\varinjlim_I X_i$“逐变量计算”。特别地，若各\Term{fiber-category}{纤维范畴}满足条件 $\mathrm L_{\Pi'_\circ}$（其中 $\Pi'_\circ$ 是一个给定\Term{cardinal}{基数}），则 $F$ 也如此。
\end{SGAresult}

9.21.10 的（容易的）详细证明留给读者；我们只指出，使用下述结果很方便，其证明是直接的。

\begin{SGAresult}{推论}{9.21.10.1}
对于 $p:E\to B$ 和 $I$，采用 9.21.10 的假设与记号。则对每个 $\alpha\in\operatorname{ob}B$，\Term{inclusion-functor}{包含函子} $E_\alpha\to E$ 与 $I$ 型\Term{inductive-limit}{归纳极限}交换。
\end{SGAresult}

回到 9.21.9 的一般条件，设 $X$ 是 $F$ 的一个\Term{object}{对象}。因为对每个 $\alpha\in\operatorname{ob}B$，$X(\alpha)$ 在 $E_\alpha$ 中\Term{accessible}{可达}，所以存在\Term{cardinal}{基数} $d\in\mathcal U$，使得对每个 $\alpha$，$X(\alpha)$ 在 $E_\alpha$ 中是 $d$-\Term{accessible}{可达}的。可以选取 $d\geq\Pi'_\circ$，从而由 9.21.10，$F$ 满足 $\mathrm L_d$。再次使用 9.21.10，在 $I$ 相对于 $d$ \Term{large-relative-to-cardinal}{足够大} 时计算 $F$ 中的\Term{inductive-limit}{归纳极限} $\varinjlim_I Y_i$，立即看出 $X$ 是 $d$-\Term{accessible}{可达}的。这就证明了 (i)。

现在分几步证明 9.21.9 (ii)（(9.21.11) 至 (9.21.16)）。设
\[
X:\alpha\longmapsto X(\alpha)
\]
是 $F$ 的一个\Term{object}{对象}。由 (9.21.4 b))，对每个 $\alpha\in\operatorname{ob}B$，可以找到一个相对于 $c$ \Term{large-relative-to-cardinal}{足够大} 的\Term{ordered-set}{有序集} $I_\alpha$，及\Term{isomorphism}{同构} $X(\alpha)=\varinjlim_{i\in I_\alpha}X(\alpha)_i$，其中 $i\mapsto X(\alpha)_i$ 是 $E_\alpha$ 中的 $I_\alpha$ 型\Term{inductive-system}{归纳系统}，各 $X(\alpha)_i$ 在 $C_\alpha$ 中。考虑有序积集 $I=\prod_\alpha I_\alpha$；显然它相对于 $c$ \Term{large-relative-to-cardinal}{足够大}，且上述\Term{inductive-system}{归纳系统}给出各 $E_\alpha$ 中的\Term{inductive-system}{归纳系统}
\[
i\longmapsto X(\alpha)_i\in\operatorname{ob}C_\alpha,
\quad i\in\operatorname{ob}I,\tag{9.21.11}
\]
它们以同一个\Term{ordered-set}{有序集} $I$ 为指标，并且对每个 $\alpha\in\operatorname{ob}B$，给出\Term{isomorphism}{同构}
\[
X(\alpha)\xrightarrow{\sim}\varinjlim_I X(\alpha)_i
\quad\text{于 }E_\alpha\text{ 中。}\tag{9.21.12}
\]
以下假定已选定数据 (9.21.11) 和 (9.21.12)。

\begin{SGAresult}{引理}{9.21.13}
在上述条件下，可以找到一个映射 $\varphi:I\to I$，使对每个 $i\in I$ 都有 $\varphi(i)\geq i$，以及一个映射 $(i,f)\mapsto\lambda(i,f)$，从 $I\times\operatorname{Fl}B$ 到 $\operatorname{Fl}E$，满足下述条件：
\SourcePage{101}
a) 对 $i\in I$ 和 $(f:\alpha\to\beta)\in\operatorname{Fl}B$，$\lambda(i,f)$ 是从 $X(\alpha)_i$ 到 $X(\beta)_{\varphi(i)}$ 的一个 $f$-\Term{morphism}{态射}，使下图交换：
\[
\begin{tikzcd}[column sep=large,row sep=large]
X(\alpha)\arrow[rr,"X(f)"] && X(\beta)\\
&&X(\beta)_{\varphi(i)}\arrow[u]\\
X(\alpha)_i\arrow[uu]\arrow[urr,"{\lambda(i,f)}"]&&\\
\alpha\arrow[rr,"f"]&&\beta .
\end{tikzcd}
\]
其中竖直\Term{arrow}{箭头}是由 (9.21.12) \Term{derivation}{导出}的\Term{canonical}{典范}\Term{morphism}{态射}。

b) 对每个 $i\in I$ 和 $(f:\alpha\to\beta)\in\operatorname{Fl}B$，下图交换：
\[
\begin{tikzcd}[column sep=large,row sep=large]
&&X(\beta)_{\varphi^2(i)}\\
X(\alpha)_{\varphi(i)}\arrow[urr,"{\lambda(\varphi(i),f)}"]&&X(\beta)_{\varphi(i)}\arrow[u]\\
X(\alpha)_i\arrow[u]\arrow[urr,"{\lambda(i,f)}"]&&\\
\alpha\arrow[rr,"f"]&&\beta .
\end{tikzcd}
\]

c) 对 $B$ 中每对相继的\Term{arrow}{箭头} $\alpha\xrightarrow{f}\beta\xrightarrow{g}\gamma$，下图交换：
\[
\begin{tikzcd}[column sep=large,row sep=large]
X(\alpha)\arrow[r,"X(f)"]&X(\beta)\arrow[r,"X(g)"]&X(\gamma)\\
&&X(\gamma)_{\varphi^2(i)}\arrow[u]\\
&X(\beta)_{\varphi(i)}\arrow[uu]\arrow[ur,"{\lambda(\varphi(i),g)}"]&X(\gamma)_{\varphi(i)}\arrow[u]\\
X(\alpha)_i\arrow[uuu]\arrow[ur,"{\lambda(i,f)}"]\arrow[urr,"{\lambda(i,gf)}"']&&\\
\alpha\arrow[r,"f"]&\beta\arrow[r,"g"]&\gamma ,
\end{tikzcd}
\]
\SourcePage{102}
其中竖直\Term{arrow}{箭头}是由 (9.21.12) \Term{derivation}{导出}的\Term{arrow}{箭头}。
\end{SGAresult}

还须注意，b) 中上方两个梯形的交换性已经包含在 a) 中，所以 b) 实际上断言所考虑图中下方三角形的交换性。

\begin{proof}[9.21.13 的证明]
设 $i\in I$，$f:\alpha\to\beta$ 是 $B$ 的一个\Term{arrow}{箭头}。考虑复合 $X(\alpha_i)\to X(\alpha)\xrightarrow{X(f)}X(\beta)\simeq\varinjlim_jX(\beta)_j$，它定义 $E$ 中的一个\Term{morphism}{态射}：
\[
X(\alpha)_i\longrightarrow f^*(X(\beta))
\simeq f^*(\varinjlim_jX(\beta)_j)
\simeq\varinjlim_jf^*(X(\beta)_j),
\]
其中最后一个等式来自 $f^*$ 是 $c$-\Term{accessible}{可达}的（9.21.3 c)）且 $I$ 相对于 $c$ \Term{large-relative-to-cardinal}{足够大}。由 (9.21.4 a))，由于 $X(\alpha)\in\operatorname{ob}C_\alpha$，所考虑的\Term{morphism}{态射}经过某个 $f^*(X(\alpha)_j)$ \Term{resolution}{分解}，其中 $j$ 先验地依赖于 $i$ 和 $f\in\operatorname{Fl}(B)$。但由 (9.21.3 b))，且 $I$ 相对于 $c$ \Term{large-relative-to-cardinal}{足够大}，可以使 $j$ 的选取与 $f$ 无关，记 $j=\varphi(i)$。$I$ 是\Term{filtered-category}{滤过的}，所以可假定 $\varphi(i)\geq i$。这样，对 $i\in I$ 和 $f\in\operatorname{Fl}F$，就找到一个\Term{morphism}{态射}
\[
X(\alpha)_i\longrightarrow f^*(X(\beta)_{\varphi(i)}),
\]
\Term{equivalence}{等价}地，找到一个 $f$-\Term{morphism}{态射}
\[
\lambda(i,f):X(\alpha)_i\longrightarrow X(\beta)_{\varphi(i)},
\]
按构造它使下图交换：
\[
\mathrm D(f)\qquad
\begin{tikzcd}[column sep=large,row sep=large]
X(\alpha)\arrow[rr,"X(f)"]&&X(\beta)\\
&&X(\beta)_{\varphi(i)}\arrow[u]\\
X(\alpha)_i\arrow[uu]\arrow[urr,"{\lambda(i,f)}"]&&\\
\alpha\arrow[rr,"f"]&&\beta .
\end{tikzcd}
\]

现在给定 $i,f,g$，考虑 9.21.13 c) 中图的下方三角形；它是否交换尚不清楚，但两个复合 $X(\alpha)_i\rightrightarrows X(\beta)_{\varphi^2(i)}$ 在与 $X(\beta)_{\varphi^2(i)}\to X(\beta)$ 复合后相等。这来自上方两个梯形及外\Term{boundary}{边界}梯形的交换性，它们分别是图 $\mathrm D(f)$、$\mathrm D(g)$ 和 $\mathrm D(gf)$。因此，由于 $X(\alpha)_{\varphi^2(i)}$ 是 $c$-\Term{accessible}{可达}的（9.21.4 a)），且 $X(\alpha)\simeq\varinjlim_{j\in I}X(\alpha)_j$，其中 $I$ 相对于 $c$ \Term{large-relative-to-cardinal}{足够大}，所以可以在 $I$ 中找到元素 $j\geq\varphi^2(i)$，使所考虑的两个\Term{arrow}{箭头}在与 $X(\beta)_{\varphi^2(i)}\to X(\beta)_j$ 复合后相等。先验地，$j$ 依赖于 $i,f,g$。但固定 $i$ 后，所有可能的 $(f,g)$ 对所成\Term{set}{集合}，
\SourcePage{103}
由 (9.21.3 a) 和 b))，其\Term{cardinal}{基数} $\geq c^2=c$；所以，由于 $I$ 相对于 $c$ \Term{large-relative-to-cardinal}{足够大}，可使 $j$ 的选取与 $f,g$ 无关，记 $j=\varphi'(i)\geq\varphi^2(i)\geq\varphi(i)$。于是对每个 $i\in I$ 和 $f:\alpha\to\beta$，令
\[
\lambda'(i,f):X(\beta)_i\longrightarrow X(\beta)_{\varphi'(i)}
\]
为复合的 $f$-\Term{morphism}{态射} $X(\alpha)_i\to X(\beta)_{\varphi(i)}\to X(\beta)_{\varphi'(i)}$，其中第二个\Term{arrow}{箭头}是\Term{transition-morphism}{过渡态射}。由构造立即可知，$(\varphi',\lambda')$ 满足对于 $(\varphi,\lambda)$ 所述的条件 9.21.13 a) 和 c)。对条件 b) 也同样进行，便知可选取 $\varphi'$，使 $(\varphi',\lambda')$ 也满足这一条件。这就完成了 (9.21.13) 的证明。
\end{proof}

\textbf{9.21.14。}现在设
\[
J\subset I
\]
是 $I$ 的一个\Term{subset}{子集}，满足下述条件：

a) $J$ 是\Term{filtered-category}{滤过的}；

b) 对每个 $j\in J$，有 $\varphi(j)\in J$；

c) 对每个 $\alpha\in\operatorname{ob}B$，$X_J(\alpha)=\varinjlim_{i\in J}X(\alpha)_i$ 在 $E_\alpha$ 中\Term{representable}{可表示}。

特别地，若 $J$ 相对于 $\Pi'_\circ$ \Term{large-relative-to-cardinal}{足够大}，则 a) 和 c) 成立（9.21.2）。由 9.21.10.1 可知，对每个 $\alpha\in\operatorname{ob}B$，$X_J(\alpha)$ 是 $E$ 中的\Term{inductive-limit}{归纳极限} $\varinjlim_{i\in J}X(\alpha)_i$。对于 $B$ 的一个\Term{arrow}{箭头} $f:\alpha\to\beta$，由 (9.21.13 b))，当 $i$ 遍历 $J$ 时，\Term{arrow}{箭头} $\lambda(i,f)$ 定义从 $(X(\alpha)_i)_{i\in J}$ 到 $(X(\beta)_i)_{i\in J}$ 的一个 \Term{ind-object-morphism}{ind-对象态射}，从而在\Term{inductive-limit}{归纳极限}上得到\Term{homomorphism}{同态}
\[
X_J(f):X_J(\alpha)\longrightarrow X_J(\beta),
\]
它显然是一个 $f$-\Term{homomorphism}{同态}。利用 (9.21.13 c))，得到传递关系
\[
X_J(gf)=X_J(g)X_J(f),
\]
所以这就定义了 $E$ 在 $B$ 上的一个\Term{section}{截面} $X_J\in\operatorname{ob}F$。最后，(9.21.13 a)) 表明各\Term{canonical-homomorphism}{典范同态}
\[
X_J(\alpha)\longrightarrow X(\alpha),\qquad\alpha\in\operatorname{ob}B,
\]
关于 $\alpha$ 是\Term{functor}{函子}性的，故有一个\Term{canonical-homomorphism}{典范同态}
\[
u_J:X_J\longrightarrow X.
\]
另一方面，若
\[
J'\supset J
\]
是 $I$ 中另一个满足上述条件 a)、b)、c) 的\Term{subset}{子集}，就得到一个\Term{canonical-homomorphism}{典范同态}
\[
u_{J',J}:X_J\longrightarrow X_{J'},
\]
\SourcePage{104}
从而各 $X_J$ 组成 $F$ 中的一个\Term{inductive-system}{归纳系统}，以 $I$ 中满足所述条件的\Term{subset}{子集} $J$ 所成的\Term{set}{集合}（按包含关系排序）为参数。最后，上述\Term{homomorphism}{同态} $u_J$ 定义从这个\Term{inductive-system}{归纳系统}到 $X$（所定义的\Term{constant-inductive-system}{常值归纳系统}）的\Term{homomorphism}{同态}。

\textbf{9.21.15。}现在设 $K$ 是由 $I$ 中满足 9.21.14 所列条件 a)、b)、c) 的\Term{subset}{子集} $J$ 组成的一个\Term{set}{集合}，并假定 $K$ 是\Term{filtered-category}{滤过的}，且其并为 $I$。利用 9.21.10.1，将问题归结为逐变量核实\Term{isomorphism}{同构}，显然有
\[
\varinjlim_{J\in K}X_J\xrightarrow{\sim}X.
\]
例如，取 $K$ 为 $I$ 中所有相对于 $\Pi'_\circ$ \Term{large-relative-to-cardinal}{足够大}、对 $\varphi$ \Term{stable}{稳定}且满足 $\operatorname{card}J\leq\Pi$ 的\Term{subset}{子集} $J$ 所成的\Term{set}{集合}。此时由 $C_\alpha^\Pi$ 的定义 (9.21.6)，对每个 $\alpha\in\operatorname{ob}B$，有 $X_J(\alpha)\in\operatorname{ob}C_\alpha^\Pi$，故 $X_J\in\operatorname{ob}F^\Pi$。因此，只要说明 $K$ 相对于 $\Pi$ \Term{large-relative-to-cardinal}{足够大}（从而是\Term{filtered-category}{滤过的}）且其并为 $I$，就证明了 9.21.9 (ii)。为此，显然只须证明 $I$ 中每个满足 $\operatorname{card}S\leq\Pi$ 的\Term{subset}{子集} $S$ 都包含于某个 $J\in K$。这由 9.21.9 (ii) 中陈述的前提假设以及下述引理得到。

\begin{SGAresult}{引理}{9.21.16}
设 $I$ 是一个\Term{ordered-set}{有序集}，$\Pi$ 是一个无限\Term{cardinal}{基数}，$\varphi:I\to I$ 是 $I$ 到自身的一个映射。

a) 假定 $I$ 是\Term{filtered-category}{滤过的}。则对于 $I$ 中每个满足 $\operatorname{card}S\leq\Pi$ 的\Term{subset}{子集} $S$，存在 $I$ 中的一个\Term{filtered-category}{滤过子集} $J\supset S$，使 $\varphi(J)\subset J$，且 $\operatorname{card}J\leq\Pi$。

b) 设 $\Pi'_\circ$ 是一个小于 $\Pi$ 的\Term{cardinal}{基数}，并假定 $I$ 相对于 $\Pi$ \Term{large-relative-to-cardinal}{足够大}。则对于 $I$ 中每个满足 $\operatorname{card}S\leq\Pi$ 的\Term{subset}{子集} $S$，存在 $I$ 中一个相对于 $\Pi'_\circ$ \Term{large-relative-to-cardinal}{足够大} 的\Term{subset}{子集} $J\supset S$，使 $\varphi(J)\subset J$，且 $\operatorname{card}J\leq\Pi$。
\end{SGAresult}
\begin{proof}
例如证明 b)（a) 的证明类似而更简单）。设 $P$ 是 $I$ 的\Term{cardinal}{基数} $\leq\Pi$ 的\Term{subset}{子集}所成的\Term{set}{集合}，$T\mapsto i_T:P\to I$ 是一个映射，使得对每个 $T\in P$，$i_T$ 是 $T$ 在 $I$ 中的一个上界；这个映射的存在只是在表达 $I$ 相对于 $\Pi$ \Term{large-relative-to-cardinal}{足够大} 的假设。设 $A$ 是一个\Term{well-ordered-set}{良序集}，满足 $\operatorname{card}A=\Pi$，且对每个 $a\in A$，所有 $a'\leq a$ 所成\Term{set}{集合}的\Term{cardinal}{基数} $<\Pi$。特别地，因为 $\Pi'_\circ<\Pi$，$A$ 相对于 $\Pi'_\circ$ \Term{large-relative-to-cardinal}{足够大}。用\Term{transfinite-recursion}{超限递归}定义映射
\[
S\longmapsto S_a:P\longrightarrow P\qquad(a\in A),
\]
所用公式如下：
\[
S_\circ=S,\qquad S_{a+1}=S_a\cup\varphi(S_a)\cup\{i_{S_a}\},
\qquad S_a=\bigcup_{a<a}S_a\quad\text{若 }a\text{ 是极限序数。}
\]
对 $S\in P$，令
\[
J(S)=\bigcup_{a\in A}S_a.
\]
\SourcePage{105}
显然，$J(S)\supset S$，$J(S)$ 对 $\varphi$ \Term{stable}{稳定}，$\operatorname{card}J(S)\leq\Pi$（因为 $\operatorname{card}A=\Pi$），最后 $J(S)$ 相对于 $\Pi'_\circ$ \Term{large-relative-to-cardinal}{足够大}（利用 $A$ 相对于 $\Pi'_\circ$ \Term{large-relative-to-cardinal}{足够大} 这一事实）。

这就证明了 9.21.16，并完成了 9.21.9 的证明。
\end{proof}

还指出 9.21.9 的一个变体。

\begin{SGAresult}{引理}{9.21.17}
记号同 9.21.9。以 $F'$ 表示 $F=\mathscr{H}om_B(B,E)$ 中由 $E$ 在 $F$ 上的\Term{cartesian-section}{笛卡尔截面}组成的\Term{full-subcategory}{满子范畴} $\mathscr{H}om_{\mathrm{cart}\,B}(B,E)$。则：

(i) $F'$ 的每个\Term{object}{对象}都\Term{accessible}{可达}。更精确地，设 $\alpha\mapsto X(\alpha)$ 是 $F'$ 的一个元素，$d$ 是一个满足 $d\geq\Pi'_\circ$、$d\geq c$ 的\Term{cardinal}{基数}，且对每个 $\alpha\in\operatorname{ob}B$，$X(\alpha)$ 都是 $E$ 的 $d$-\Term{accessible-object}{可达对象}。则 $X$ 是 $F'$ 的 $d$-\Term{accessible-object}{可达对象}。

(ii) 假定 $\Pi\leq c$，且\Term{functor}{函子} $f^*:E_\beta\to E_\alpha$ 都是 $\Pi'_\circ$-\Term{accessible}{可达}的；或者，各\Term{category}{范畴} $E_\alpha$ 对小\Term{filtered-inductive-limit}{滤过归纳极限}\Term{stable}{稳定}，且\Term{functor}{函子} $f^*:E_\beta\to E_\alpha$ 与这些\Term{inductive-limit}{归纳极限}交换。则 $F'$ 的每个\Term{object}{对象} $X$ 都\Term{isomorphism}{同构}于形如 $\varinjlim_I X_i$ 的\Term{object}{对象}，其中对每个 $i\in I$，$X_i$ 是 $F'$ 的\Term{full-subcategory}{满子范畴} $F'^\Pi=F'\cap F^{\Pi_\circ}$ 的一个\Term{object}{对象}，且 $I$ 相对于 $\Pi$ \Term{large-relative-to-cardinal}{足够大}。
\end{SGAresult}

证明与 9.21.9 完全类似，我们只指出需要修改后者的地方。先注意：

\begin{SGAresult}{引理}{9.21.18}
在 9.21.10 的陈述中，将 $F=\mathscr{H}om_B(B,E)$ 换成 $F'=\mathscr{H}om_{\mathrm{cart}\,B}(B,E)$ 后，该陈述仍然成立，只要假定\Term{inverse-image-functor}{逆像函子} $f^*:E_\beta\to E_\alpha$ 与 $I$ 型\Term{inductive-limit}{归纳极限}交换。
\end{SGAresult}

所以在 9.21.17 的条件下，若 $d$ 是满足 $d\geq c$、$d\geq\Pi'_\circ$ 的一个\Term{cardinal}{基数}，则由上文及假设 (9.21.3 c))，对每个相对于 $d$ \Term{large-relative-to-cardinal}{足够大} 的\Term{ordered-set}{有序集} $I$，$I$ 型\Term{inductive-limit}{归纳极限}在 $F'$ 中\Term{representable}{可表示}，并逐变量计算。于是用一个直接的论证便得结论 9.21.17 (i)。

为证明 (ii)，必须补充 9.21.13。

\begin{SGAresult}{引理}{9.21.19}
在 9.21.13 的条件下，假定 $X\in\operatorname{ob}F'$，即 $X$ 是 $E$ 在 $B$ 上的一个\Term{cartesian-section}{笛卡尔截面}。则可以加强结论，断言存在一个函数 $\mu$，它对每个 $i\in I$ 及 $B$ 的每个\Term{arrow}{箭头} $f:\alpha\to\beta$，关联 $E_\alpha$ 中的一个\Term{morphism}{态射} $\mu(i,f):f^*X(\beta)_i\to X(\alpha)_{\varphi(i)}$，使得：

d) 对每个 $i\in I$ 及 $B$ 的每个\Term{arrow}{箭头} $f:\alpha\to\beta$，下述两个图交换：
\SourcePage{106}
\[
\begin{tikzcd}[column sep=large,row sep=large]
&f^*X(\beta)_{\varphi^2(i)}\\
X(\alpha)_{\varphi(i)}\arrow[ur,"{\lambda(\varphi(i),f)}"]&\\
&f^*X(\beta)_i\arrow[uu]\arrow[ul,"{\mu(i,f)}"]
\end{tikzcd}
\qquad
\begin{tikzcd}[column sep=large,row sep=large]
X(\alpha)_{\varphi^2(i)}&\\
&f^*X(\alpha)_{\varphi(i)}\arrow[ul,"{\mu(\varphi(i),f)}"']\\
x(\alpha)_i\arrow[uu]\arrow[ur,"{\lambda(i,j)}"']& .
\end{tikzcd}
\]
\end{SGAresult}
\begin{proof}
为证明这一点，如 9.21.3 中那样进行，考虑复合\Term{morphism}{态射}
\[
f^*(X(\beta)_i)\longrightarrow f^*(X(\beta))
\xrightarrow{X(f)^{-1}}X(\alpha),
\]
其中（如同在写出条件 d) 时已经做的那样）在记号上将 $E$ 中的一个 $f$-\Term{morphism}{态射} $R\to S$ 与 $E_\alpha$ 中\Term{correspondence}{对应}的\Term{morphism}{态射} $R\to f^*(S)$ 认同。由于 $X(\alpha)=\varinjlim_jX(\alpha)_j$，且 $I$ 相对于 $c$ \Term{large-relative-to-cardinal}{足够大}，可以使上述\Term{morphism}{态射}经过某个 $X(\alpha)_j$ \Term{resolution}{分解}，其中 $j$ 先验地依赖于 $i$ 和 $f$，但可选取为与 $f$ 无关，记为 $\psi(i)$。必要时增大 9.21.13 中的函数 $\varphi$，可假定 $\psi=\varphi$。仿照对 9.21.13 的条件 b)、c) 的处理，必要时进一步增大映射 $\varphi$，就可假定 9.21.19 d) 的两个图交换。
\end{proof}

\textbf{9.21.20。}如 9.21.19 那样选取映射 $\varphi$ 后，重新采用 9.21.14 的论证，但现在必须假定 $J$ 除了满足所列条件 a)、b)、c) 外，还满足：

d) 对 $B$ 中每个\Term{arrow}{箭头} $f:\alpha\to\beta$，\Term{functor}{函子} $f^*:E_\beta\to E_\alpha$ 与 $J$ 型\Term{inductive-limit}{归纳极限}交换。

断言，在这个附加条件下，$E$ 在 $B$ 上的\Term{section}{截面} $X_J$ 是\Term{cartesian}{笛卡尔}的，即对于 $B$ 的每个\Term{arrow}{箭头} $f:\alpha\to\beta$，\Term{morphism}{态射}
\[
X_J(\alpha)\longrightarrow f^*X_J(\beta)\tag{*}
\]
是\Term{isomorphism}{同构}。事实上，由上述条件 d)，所考虑\Term{morphism}{态射}的目标认同于 $\varinjlim_{i\in J}f^*(X(\beta)_i)$，而该\Term{morphism}{态射}由 $E_\alpha$ 中的 \Term{ind-object-morphism}{ind-对象态射}
\[
(X(\alpha)_i)_{i\in J}\longrightarrow (f^*)(X(\beta)_i)_{i\in J}
\]
取 $J$ 上的\Term{inductive-limit}{归纳极限}而得，后者来自 $i\in J$ \Term{correspondence}{对应}的\Term{morphism}{态射} $\lambda(i,f):X(\alpha)_i\to f^*(X(\beta)_{\varphi(i)})$。9.21.19 中明确给出的条件保证，上述 \Term{ind-object}{ind-对象}\Term{homomorphism}{同态}实际上是\Term{isomorphism}{同构}，其逆由 $i\in J$ \Term{correspondence}{对应}的 $\mu(i,f)$ 系统所\Term{derivation}{导出}的\Term{homomorphism}{同态}给出。这就证明了 (*) 是\Term{isomorphism}{同构}的断言，故 $X\in\operatorname{ob}F'$。

\SourcePage{107}
\textbf{9.21.21。}现在假定\Term{functor}{函子} $f^*:E_\beta\to E_\alpha$ 是 $\Pi'_\circ$-\Term{accessible}{可达}的。于是，若 $J$ 相对于 $\Pi'_\circ$ \Term{large-relative-to-cardinal}{足够大}、对 $\varphi$ \Term{stable}{稳定}、且 $\operatorname{card}J\leq\Pi$，则 9.21.20 中考虑的关于 $J$ 的条件成立，于是如 9.21.15 中那样完成 9.21.17 的证明。

\begin{SGAresult}{定理}{9.22}
设 $p:E\to B$ 是一个\Term{fibred-functor}{纤维化函子}，其中 $B$ 是一个本质上\Term{equivalence}{等价}于\Term{small-category}{小范畴}的\Term{category}{范畴}，且对 $B$ 的每个\Term{arrow}{箭头} $f:\alpha\to\beta$，\Term{inverse-image-functor}{逆像函子} $f^*:E_\beta\to E_\alpha$ 都\Term{accessible}{可达}（9.2）。进一步假定，对每个 $\alpha\in\operatorname{ob}E$，\Term{fiber-category}{纤维范畴} $E_\alpha$ 具有\Term{cardinal-filtration}{基数滤过}（9.12），且 $E_\alpha$ 的每个\Term{object}{对象}都\Term{accessible}{可达}（9.3）。则\Term{category}{范畴} $F=\mathscr{H}om_B(B,E)$ 和 $F'=\mathscr{H}om_{\mathrm{cart}\,B}(B,E)$ 都具有\Term{generating-by-strict-epimorphisms}{由严格满态射生成}的小\Term{full-subcategory}{满子范畴}（7.1），且其每个\Term{object}{对象}都\Term{accessible}{可达}。
\end{SGAresult}
\begin{proof}
显然，若将 $B$ 换成一个使\Term{inclusion-functor}{包含函子} $B_\circ\to B$ 为\Term{equivalence}{等价}的\Term{full-subcategory}{满子范畴} $B_\circ$，并将 $E$ 换成 $E_\circ=E\times_B B_\circ$，则陈述本质上不变。因此可以假定 $B$ 是\Term{small-category}{小范畴}。

对每个 $\alpha\in\operatorname{ob}B$，设 $(\operatorname{Filt}^\Pi(E))_{\Pi\geq\Pi_\alpha}$ 是 $E_\alpha$ 的一个\Term{cardinal-filtration}{基数滤过}。设 $c_\circ\in\mathcal U$ 是一个满足下述条件的\Term{cardinal}{基数}：
\[
\left\{\begin{aligned}
{}&c_\circ\geq\sup_{\alpha\in\operatorname{ob}B}\Pi_\alpha,
\quad c_\circ\geq\operatorname{card}\operatorname{Fl}B;\\
&\text{对 }B\text{ 的每个箭头 }f:\alpha\to\beta,
\quad f^*:E_\beta\to E_\alpha\text{ 是 }c_\circ\text{-可达的；}\\
&\text{对每个 }\alpha\in\operatorname{ob}B,
\quad\operatorname{Filt}^{\Pi_\alpha}(E_\alpha)\text{ 的对象都是 }c_\circ\text{-可达的。}
\end{aligned}\right.\tag{9.22.1}
\]
令
\[
c=2^{c_\circ},
\]
从而 $c>c_\circ$，并对每个 $\alpha\in\operatorname{ob}B$，令
\[
C_\alpha=\operatorname{Filt}^c(E_\alpha).
\]
断言，取 $\Pi=\Pi_\circ=c$，则 9.21.9 和 9.21.17 的前提条件都满足。对 (9.21.2) 和 (9.21.3)，这是显然的。(9.21.4 a)) 由 9.17 得到，(9.21.4 b)) 则由\Term{cardinal-filtration}{基数滤过}的条件 9.12 c) 得到。此外注意，\Term{cardinal-filtration}{基数滤过}的条件 9.12 b) 蕴涵，对每个 $\alpha\in\operatorname{ob}B$，有 $C_\alpha^\Pi=C_\alpha=\operatorname{Filt}^c(E_\alpha)$，其中 $C_\alpha^\Pi$ 由 (9.21.6) 定义。于是 9.22 由 9.21.9 和 9.21.17 得到。更精确地，我们已经证明下述推论。
\end{proof}

\begin{SGAresult}{推论}{9.23}
在 9.22 的条件下，若 $c_\circ\in\mathcal U$ 是一个满足条件 (9.22.1) 的\Term{cardinal}{基数}，且 $c=2^{c_\circ}$，则 $F$（分别为 $F'$）中由对每个 $\alpha\in\operatorname{ob}B$ 都满足 $X(\alpha)\in\operatorname{ob}\operatorname{Filt}^c(E)$ 的 $X$ 组成的\Term{subcategory}{子范畴} $F^c$（分别为 $F'^c$）是\Term{essentially-small-category}{本质小}的，且\Term{generating-by-strict-epimorphisms}{由严格满态射生成}；更精确地，$F$（分别为 $F'$）的每个\Term{object}{对象} $X$ 都\Term{isomorphism}{同构}于形如 $\varinjlim_I X_i$ 的\Term{object}{对象}，其中各 $X_i$ 在 $F^c$（分别在 $F'^c$）中，且 $I$ 是一个相对于 $c$ \Term{large-relative-to-cardinal}{足够大} 的\Term{ordered-set}{有序集}。
\end{SGAresult}

\SourcePage{108}
若稍微加强 9.22 的假设，还可以大幅度精化 9.22 和 9.23。

\begin{SGAresult}{推论}{9.24}
在 9.22 的条件下，假定每个\Term{category}{范畴} $E_\alpha$（$\alpha\in\operatorname{ob}E$）都对小\Term{filtered-inductive-limit}{滤过归纳极限}\Term{stable}{稳定}，而且对每个 $\Pi\geq\Pi_\alpha$，$\operatorname{Filt}^\Pi(E_\alpha)$ 都对以满足 $\operatorname{card}I\leq\Pi$ 的\Term{filtered-category}{滤过有序集} $I$ 为指标的\Term{filtered-inductive-limit}{滤过归纳极限}\Term{stable}{稳定}（这略微加强了\Term{cardinal-filtration}{基数滤过}的条件 9.12 b)）。在所考虑的是 $F'$ 的情形下，进一步假定对 $B$ 的每个\Term{arrow}{箭头} $f:\alpha\to\beta$，\Term{functor}{函子} $f^*:E_\beta\to E_\alpha$ 与小\Term{filtered-inductive-limit}{滤过归纳极限}交换。最后设 $c_\circ\in\mathcal U$ 是一个满足条件 (9.22.1) 的\Term{cardinal}{基数}，并对每个\Term{cardinal}{基数} $\Pi\geq c=2^{c_\circ}$，考虑 $F$（分别为 $F'$）中由对每个 $\alpha\in\operatorname{ob}B$ 都满足 $X(\alpha)\in\operatorname{Filt}^\Pi(E_\alpha)$ 的 $X$ 组成的\Term{strictly-full-subcategory}{严格满子范畴} $F^\Pi$（分别为 $F'^\Pi$）。则这些\Term{subcategory}{子范畴}定义 $F$（分别为 $F'$）的一个\Term{cardinal-filtration}{基数滤过}（9.12）。
\end{SGAresult}
\begin{proof}
必须核实 9.12 的条件 a)、b)、c)。条件 a)、b) 来自各 $E_\alpha$ 的给定\Term{cardinal-filtration}{基数滤过}的类似条件，以及 9.21.10（分别为 9.21.18，并考虑到 (9.22.1) 的第二个条件）。剩下证明 c)，其第一部分在取 $\Pi'_\circ=0$、$\Pi_\circ=c$ 后立即由 9.21.9（分别为 9.21.17）得到。还须证明，若有两个\Term{cardinal}{基数} $\Pi'\geq\Pi\geq c$，则对 $F^{\Pi'}$（分别为 $F'^{\Pi'}$）中的每个 $X$，有 $X\simeq\varinjlim_{i\in K}X_i$，其中各 $X_i$ 在 $F^\Pi$（分别在 $F'^\Pi$）中，$K$ 相对于 $\Pi$ \Term{large-relative-to-cardinal}{足够大}，且 $\operatorname{card}K\leq\Pi'^\Pi$。

为此，重新采用 9.21.9 (ii) 与 9.21.17 (ii) 的证明。由 $E_\alpha$ 的\Term{cardinal-filtration}{基数滤过}的条件 9.12 c)，可以假定每个 $I_\alpha$ 的\Term{cardinal}{基数} $\leq\Pi'$，所以有
\[
\operatorname{card}I\leq(\Pi'^\Pi)^{\operatorname{card}\operatorname{ob}B}
\leq(\Pi'^\Pi)^c=\Pi'^{\Pi c}=\Pi'^\Pi.
\]
9.21.15、分别 9.21.21 中使用的指标集 $K$ 是 $I$ 中所有满足下述条件的\Term{subset}{子集} $J$ 所成的\Term{set}{集合}：$J$ 是\Term{filtered-category}{滤过的}（即相对于 $\Pi'_\circ=0$ \Term{large-relative-to-cardinal}{足够大}）、对 $\varphi$ \Term{stable}{稳定}、且 $\operatorname{card}J\leq\Pi$。立即有
\[
\operatorname{card}K\leq(\operatorname{card}I)^\Pi
\leq(\Pi'^\Pi)^\Pi=\Pi'^{\Pi^2}=\Pi'^\Pi,
\]
这就完成了 9.24 的证明。
\end{proof}

最后，宜给出一个摆脱了 9.22 中略具技术性的假设的陈述：将它们换成各 $E_\alpha$ 同时满足 9.11 中的假设（它们保证各 $E$ 的\Term{object}{对象}的\Term{accessible}{可达}性）与 9.13 中的假设（它们保证 $E$ 中存在\Term{cardinal-filtration}{基数滤过}）。

\begin{SGAresult}{推论}{9.25}
设 $p:E\to B$ 是一个\Term{fibred-functor}{纤维化函子}，其中 $B$ 是一个\Term{equivalence}{等价}于\Term{small-category}{小范畴}的\Term{category}{范畴}，且对 $B$ 的每个\Term{arrow}{箭头} $f:\alpha\to\beta$，\Term{functor}{函子} $f^*:E_\beta\to E_\alpha$ \Term{accessible}{可达}（9.2）。进一步假定，对每个 $\alpha\in\operatorname{ob}B$，\Term{fiber-category}{纤维范畴} $E_\alpha$ 满足下述条件：
\SourcePage{109}
a) $E_\alpha$ 具有一个\Term{generating-by-strict-epimorphisms}{由严格满态射生成}的小\Term{full-subcategory}{满子范畴}（7.1）。

b) $E_\alpha$ 对\Term{fiber-product}{纤维积}、平行双\Term{arrow}{箭头}的\Term{equalizer}{等化子}、两个\Term{object}{对象}的和、平行双\Term{arrow}{箭头}的\Term{coequalizer}{余等化子}，以及小\Term{filtered-inductive-limit}{滤过归纳极限}\Term{stable}{稳定}。\footnote[3]{最后这个条件很可能没有必要，参见 9.13.2。}

c) $E_\alpha$ 的平行双\Term{arrow}{箭头}\Term{category}{范畴}上的\Term{functor}{函子} $\operatorname{Ker}(u,v)$ \Term{accessible}{可达}（例如，它与小\Term{filtered-inductive-limit}{滤过归纳极限}交换）。

d) $E_\alpha$ 的每个\Term{strict-epimorphism}{严格满态射}（10.2）都是\Term{universally-strict-epimorphism}{普适严格满态射}。

在这些条件下，\Term{category}{范畴} $F=\mathscr{H}om_B(B,E)$ 与 $F'=\mathscr{H}om_{\mathrm{cart}\,B}(B,E)$ 都具有\Term{generating-by-strict-epimorphisms}{由严格满态射生成}的小\Term{subcategory}{子范畴}，且它们的所有\Term{object}{对象}都\Term{accessible}{可达}（9.3）。
\end{SGAresult}

此外，9.24 给出：

\begin{SGAresult}{推论}{9.26}
在 9.25 的条件下，并在考虑 $F'=\mathscr{H}om_{\mathrm{cart}\,B}(B,E)$ 的情形下，假定\Term{functor}{函子} $f^*:E_\beta\to E_\alpha$ 与小\Term{filtered-inductive-limit}{滤过归纳极限}交换。则 $F$（分别为 $F'$）具有一个\Term{cardinal-filtration}{基数滤过}，它可以用 9.24 的过程明确给出。
\end{SGAresult}

\SGAsection{10}{术语表}
为方便读者，这里汇集前面各款中使用的若干术语的定义。以 $C$ 表示一个\Term{category}{范畴}。

\textbf{10.1。\Term{cartesian}{笛卡尔}、\Term{cocartesian}{余笛卡尔}。} $C$ 中的图
\[
\begin{tikzcd}
A\arrow[r]\arrow[d]&B\arrow[d]\\
C\arrow[r]&D
\end{tikzcd}
\]
称为\Term{cartesian}{笛卡尔}的，若它交换，且从 $A$ 到\Term{fiber-product}{纤维积} $C\times_D B$ 的\Term{canonical}{典范}\Term{morphism}{态射}是\Term{isomorphism}{同构}。若 $A$ 的\Term{opposite-category}{对偶范畴}中\Term{correspondence}{对应}的图是\Term{cartesian}{笛卡尔}的，则称它为\Term{cocartesian}{余笛卡尔}的。

\textbf{10.2。\Term{epimorphism}{满态射}、\Term{strict-epimorphism}{严格满态射}，等等。}参见 10.3。

\SourcePage{110}
\textbf{10.3。\Term{epimorphic-family}{满态射族}、\Term{strict-epimorphic-family}{严格满态射族}，等等。} $C$ 中的一个有共同目标的\Term{arrow}{箭头}族
\[
f_i:A_i\longrightarrow B,\qquad i\in I,
\]
称为\Term{epimorphic-family}{满态射族}，若对 $C$ 的每个\Term{object}{对象} $C$，由各 $f_i$ 诱导的映射
\[
\operatorname{Hom}_{\mathcal A}(B,C)
\longrightarrow\prod_i\operatorname{Hom}_{\mathcal A}(A_i,C)\tag{10.3.1}
\]
是\Term{injection}{单射}。

若映射 (10.3.1) 的像由满足下述条件的族 $(g_i)$ 组成，就称所考虑的族为\Term{strict-epimorphic-family}{严格满态射族}：对 $C$ 的每个\Term{object}{对象} $R$、每对指标 $i,j\in I$，以及每对满足 $f_iu=f_jv$ 的\Term{arrow}{箭头} $u:R\to A_i$、$v:R\to A_j$，也有 $g_iu=g_jv$。若对 $i,j\in I$，\Term{fiber-product}{纤维积} $A_i\times_BA_j$ 都\Term{representable}{可表示}，则\Term{equivalence}{等价}于说，(10.3.1) 的\Term{essential-image}{本质像}由对每对指标 $i,j\in I$ 都满足 $g_i\operatorname{pr}_1=g_i\operatorname{pr}_2$ 的 $(g_i)$ 组成，其中 $\operatorname{pr}_1,\operatorname{pr}_2$ 是 $A_i\times_B A_j$ 的两个\Term{projective-module}{投射}。若上述条件成立，即该族是\Term{strict-epimorphic-family}{严格满态射族}，并且各\Term{fiber-product}{纤维积} $A_i\times_A A_j$ 都\Term{representable}{可表示}，就称族 $(f_i)_{i\in I}$ 为\Term{effective-epimorphic-family}{有效满态射族}。

若各\Term{morphism}{态射} $f_i$ \Term{quarrable-morphism}{可作基变换}（10.3），且对每个\Term{arrow}{箭头} $B'\to B$，由族 $(f_i)_{i\in I}$ 作\Term{base-change}{基变换} $B'\to B$ 得到的\Term{arrow}{箭头}族 $f'_i:A'_i\to A_i\times_BB'$、$i\in I$ 是\Term{epimorphic-family}{满态射族}（分别为\Term{effective-epimorphic-family}{有效满态射族}），则称族 $(f_i)_{i\in I}$ 为\Term{universally-epimorphic-family}{普适满态射族}（分别为\Term{universally-effective-epimorphic-family}{普适有效满态射族}）。

\Term{universally-strict-epimorphic-family}{普适严格满态射族}的概念将在第 II 讲 2.5、2.6 中定义。注意，当各\Term{morphism}{态射} $f_i$ \Term{quarrable-morphism}{可作基变换}时（例如 $C$ 中的\Term{fiber-product}{纤维积}\Term{representable}{可表示}时），这一概念与\Term{universally-effective-epimorphic-family}{普适有效满态射族}的概念相同（利用第 II 讲 2.4）。在一般情形下，所考虑的\Term{epimorphism}{满态射}性概念的六个变体之间，有下述逻辑蕴涵关系：
\[
\resizebox{\linewidth}{!}{\begin{tikzcd}[ampersand replacement=\&,column sep=small,row sep=large]
\&\&\text{普适有效满态射族}\arrow[dll,Rightarrow]\arrow[drr,Rightarrow]\&\&\\
\text{有效满态射族}\arrow[ddrr,Rightarrow]\&\&\&\&\text{普适严格满态射族}\arrow[dl,Rightarrow]\arrow[dr,Rightarrow]\\
\&\&\&\text{严格满态射族}\arrow[dl,Rightarrow]\&\&\text{普适满态射族}\arrow[dlll,Rightarrow]\\
\&\&\text{满态射族}\&\&\&
\end{tikzcd}}
\]

\SourcePage{111}
$\mathcal A$ 中的一个\Term{morphism}{态射} $f:A\to B$ 称为\Term{epimorphism}{满态射}（分别为\Term{strict-epimorphism}{严格满态射}、\Term{effective-epimorphism}{有效满态射}、\Term{universal-epimorphism}{普适满态射}、\Term{universally-effective-epimorphism}{普适有效满态射}、\Term{universally-strict-epimorphism}{普适严格满态射}），若仅由 $f$ 一个元素组成的\Term{morphism}{态射}族是\Term{epimorphic-family}{满态射族}（分别为……）。

\textbf{10.4。\Term{monomorphic-family}{单态射族}、\Term{strict-monomorphic-family}{严格单态射族}，等等。} $C$ 中的一个\Term{arrow}{箭头}族 $f_i:A_i\to B$ 称为\Term{monomorphic-family}{单态射族}（分别为\Term{strict-monomorphic-family}{严格单态射族}，……），若将它看作\Term{opposite-category}{对偶范畴} $C^\circ$ 的\Term{arrow}{箭头}族时，它是\Term{epimorphic-family}{满态射族}（分别为\Term{strict-epimorphic-family}{严格满态射族}，……）。$\mathcal A$ 的一个\Term{arrow}{箭头} $f:A\to B$ 称为\Term{monomorphism}{单态射}（分别为\Term{strict-monomorphism}{严格单态射}，……），若仅由 $f$ 组成的族是\Term{monomorphic-family}{单态射族}（分别为\Term{strict-monomorphic-family}{严格单态射族}，……）；即 $f$ 作为\Term{opposite-category}{对偶范畴} $C^\circ$ 的一个\Term{arrow}{箭头}，是\Term{epimorphism}{满态射}（分别为\Term{strict-epimorphism}{严格满态射}，……）。

\textbf{10.5。\Term{monomorphism}{单态射}、\Term{strict-monomorphism}{严格单态射}，等等。}参见 10.4。

\textbf{10.6。\Term{projector}{投影子}、\Term{projector}{投影子}的像、\Term{object}{对象}的\Term{direct-factor}{直因子}。}设 $f:X\to X$ 是 $C$ 的一个\Term{object}{对象}的\Term{endomorphism}{自同态}。若 $f^2=f$，就称 $f$ 是\Term{projector}{投影子}。此时，对 $(f,\operatorname{id}_X)$ 有\Term{representable}{可表示}的\Term{coequalizer}{余等化子} $X'$ 的充要条件是它有\Term{representable}{可表示}的\Term{equalizer}{等化子} $X''$；在此情形下，$C$ 中有唯一的\Term{isomorphism}{同构} $u:X'\to X''$，使 $iup=f$，其中 $p:X\to X'$ 和 $i:X''\to X$ 是\Term{canonical}{典范}\Term{morphism}{态射}。通常把 $X'$ 和 $X''$ 认同，并称为\Term{projector}{投影子} $f$ 的像；若 $\operatorname{Ker}(f,\operatorname{id}_X)$ \Term{representable}{可表示}，即 $\operatorname{Coker}(f,\operatorname{id}_X)$ \Term{representable}{可表示}，就说\Term{projector}{投影子} $f$ 有像。$X$ 的一个\Term{subobject}{子对象}（分别为\Term{quotient-object}{商对象}）$Y$ 称为 $X$ 的\Term{direct-factor}{直因子}，若可以找到 $X$ 上的一个\Term{projector}{投影子} $f$，使它经过 $Y$ \Term{resolution}{分解}，并且以 $Y$ 作为它的像，视为 $X$ 的\Term{subobject}{子对象}（分别视为\Term{quotient-object}{商对象}）。有时也滥用术语，把 $C$ 的一个\Term{object}{对象}称为 $X$ 的\Term{direct-factor}{直因子}，若它\Term{isomorphism}{同构}于 $X$ 上某个\Term{projector}{投影子}的像。

\textbf{10.7。\Term{quarrable-morphism}{可作基变换}。} $C$ 的一个\Term{arrow}{箭头} $f:A\to B$ 称为\Term{quarrable-morphism}{可作基变换}的，若对 $C$ 的每个\Term{arrow}{箭头} $B'\to B$，\Term{fiber-product}{纤维积} $A\times_BB'$ 在 $C$ 中\Term{representable}{可表示}。

\textbf{10.8。\Term{quotient-object}{商对象}、\Term{strict-quotient}{严格商对象}，等等。}设 $A$ 是 $C$ 的一个\Term{object}{对象}。两个以 $A$ 为源的\Term{epimorphism}{满态射} $f:A\to B$、$f':A\to B'$ 称为\Term{equivalence}{等价}的，若存在\Term{isomorphism}{同构} $u:B\xrightarrow{\sim}B'$，使 $uf=f'$。这在以 $A$ 为源的\Term{epimorphism}{满态射}集上给出一个\Term{equivalence-relation}{等价关系}，其\Term{equivalence-class}{等价类}称为 $A$ 的商，或\Term{quotient-object}{商对象}。对于 $A$ 的每个\Term{quotient-object}{商对象}，通常假定已选定该\Term{equivalence-class}{等价类}的一个元素，设为 $f:A\to B$，并且经常（滥用术语）称 $B$ 为 $A$ 的\Term{quotient-object}{商对象}（或者称 $f:A\to B$ 为 $A$ 的\Term{quotient-object}{商对象}）。若\Term{morphism}{态射} $f:A\to B$ 是\Term{strict-epimorphism}{严格满态射}（分别为\Term{effective-epimorphism}{有效满态射}、\Term{universal-epimorphism}{普适满态射}，……）（10.3），就称 $B$ 为\Term{strict-quotient}{严格商对象}（分别为\Term{effective-quotient}{有效商对象}、\Term{universal-quotient}{普适商对象}，……）。

\SourcePage{112}
\textbf{10.9。\Term{equivalence-relation}{等价关系}。}设 $F,G$ 是 $C$ 上的两个\Term{set}{集合}\Term{presheaf}{预层}。图 $p_1,p_2:F\rightrightarrows G$ 称为 $G$ 上的一个\Term{equivalence-relation}{等价关系}，若对 $C$ 的每个\Term{object}{对象} $A$，映射
\[
(p_1(A),p_2(A)):F(A)\longrightarrow G(A)\times G(A)
\]
诱导从 $F(A)$ 到\Term{set}{集合} $G(A)$ 上某个\Term{equivalence-relation}{等价关系}的图的\Term{bijection}{双射}。

$C$ 中的一个图 $p_1,p_2:B\rightrightarrows C$ 称为 $C$ 上的一个\Term{equivalence-relation}{等价关系}，若\Term{correspondence}{对应}的\Term{presheaf}{预层}图是\Term{equivalence-relation}{等价关系}。若 $C$ 中的\Term{finite-product}{有限积}和\Term{fiber-product}{纤维积}都\Term{representable}{可表示}，则与\Term{finite-product}{有限积}和\Term{fiber-product}{纤维积}交换的\Term{functor}{函子} $\Phi:C\to C'$，把 $C$ 上的\Term{equivalence-relation}{等价关系}变成 $\Phi(C)$ 上的\Term{equivalence-relation}{等价关系}。

设 $\Pi:C\to D$ 是 $\mathcal A$ 的一个\Term{morphism}{态射}，且\Term{fiber-product}{纤维积} $C\times_D C$ \Term{representable}{可表示}。则图 $\operatorname{pr}_1,\operatorname{pr}_2:C\times_D C\rightrightarrows C$ 是 $C$ 上的一个\Term{equivalence-relation}{等价关系}。

\textbf{10.10。\Term{effective-equivalence-relation}{有效等价关系}、\Term{universally-effective}{普适有效}\Term{equivalence-relation}{等价关系}。} $C$ 的一个\Term{object}{对象} $A$ 上的\Term{equivalence-relation}{等价关系} $p_1,p_2:R\rightrightarrows A$ 称为\Term{effective-equivalence-relation}{有效等价关系}，若存在\Term{morphism}{态射} $\Pi:A\to B$，使得方块
\[
\begin{tikzcd}
R\arrow[r,"p_2"]\arrow[d,"p_1"']&A\arrow[d,"\pi"]\\
A\arrow[r]&B
\end{tikzcd}
\]
既是\Term{cartesian}{笛卡尔}的，又是\Term{cocartesian}{余笛卡尔}的。（\Term{morphism}{态射} $\pi$ 是对 $(p_1,p_2)$ 的\Term{coequalizer}{余等化子}，因此在相差唯一\Term{isomorphism}{同构}的意义下确定。）\Term{morphism}{态射} $\pi$ 此时是\Term{effective-epimorphism}{有效满态射}；若它是\Term{universally-effective-epimorphism}{普适有效满态射}（10.3），则称该\Term{equivalence-relation}{等价关系}为\Term{universally-effective}{普适有效}的。

\textbf{10.11。\Term{subobject}{子对象}、\Term{strict-subobject}{严格子对象}，等等。}这些概念分别是\Term{quotient-object}{商对象}、\Term{strict-quotient}{严格商对象}等概念（10.8）的\Term{duality}{对偶}。

\bigskip
\begin{center}\textbf{参考文献}\end{center}
\noindent [1] B.~Mitchell，\emph{\Term{category}{范畴}论}（Theory of categories），Academie Press（1965）。
